\documentclass[12pt,a4paper]{ctexart}

% === 核心包 ===
\usepackage{geometry}
\usepackage{graphicx}
\usepackage{booktabs}
\usepackage{longtable}
\usepackage{array}
\usepackage{calc}
\usepackage{hyperref}
\usepackage{xcolor}
\usepackage{enumitem}
\usepackage{etoolbox}
% longtable 用小号字
\AtBeginEnvironment{longtable}{\small}
% pandoc 输出的列表紧排命令
\providecommand{\tightlist}{%
  \setlength{\itemsep}{0pt}\setlength{\parskip}{0pt}}
% 特殊 Unicode 符号映射
\usepackage{newunicodechar}
\usepackage{amssymb}
% 希腊字母
\newunicodechar{α}{\ensuremath{\alpha}}
\newunicodechar{β}{\ensuremath{\beta}}
\newunicodechar{κ}{\ensuremath{\kappa}}
\newunicodechar{σ}{\ensuremath{\sigma}}
\newunicodechar{χ}{\ensuremath{\chi}}
% 运算符
\newunicodechar{⭐}{\ensuremath{\bigstar}}
\newunicodechar{≤}{\ensuremath{\leq}}
\newunicodechar{≥}{\ensuremath{\geq}}
\newunicodechar{≈}{\ensuremath{\approx}}
\newunicodechar{≠}{\ensuremath{\neq}}
% 下标
\newunicodechar{₀}{\textsubscript{0}}\newunicodechar{₁}{\textsubscript{1}}
\newunicodechar{₂}{\textsubscript{2}}\newunicodechar{₃}{\textsubscript{3}}
\newunicodechar{₄}{\textsubscript{4}}\newunicodechar{₈}{\textsubscript{8}}
\newunicodechar{₆}{\textsubscript{6}}\newunicodechar{₇}{\textsubscript{7}}
\newunicodechar{ₑ}{\textsubscript{e}}\newunicodechar{ₗ}{\textsubscript{l}}
\newunicodechar{ₓ}{\textsubscript{x}}
% 上标
\newunicodechar{⁺}{\ensuremath{^{+}}}
\newunicodechar{⁻}{\ensuremath{^{-}}}
\newunicodechar{⁴}{\textsuperscript{4}}
\newunicodechar{⁶}{\textsuperscript{6}}
\newunicodechar{⁷}{\textsuperscript{7}}
\newunicodechar{ᵧ}{\textsuperscript{$\gamma$}}
% 带圈数字 ①–⑧
\newunicodechar{①}{\textcircled{\raisebox{-0.5pt}{1}}}
\newunicodechar{②}{\textcircled{\raisebox{-0.5pt}{2}}}
\newunicodechar{③}{\textcircled{\raisebox{-0.5pt}{3}}}
\newunicodechar{④}{\textcircled{\raisebox{-0.5pt}{4}}}
\newunicodechar{⑤}{\textcircled{\raisebox{-0.5pt}{5}}}
\newunicodechar{⑥}{\textcircled{\raisebox{-0.5pt}{6}}}
\newunicodechar{⑦}{\textcircled{\raisebox{-0.5pt}{7}}}
\newunicodechar{⑧}{\textcircled{\raisebox{-0.5pt}{8}}}
% 状态图标
\newunicodechar{✅}{\checkmark}
\newunicodechar{⏳}{[TBD]}
\newunicodechar{⚠}{\textbf{/!\\}}
% 罗马数字
\newunicodechar{Ⅱ}{II}
% 补充数学/特殊符号
\newunicodechar{²}{\textsuperscript{2}}
\newunicodechar{³}{\textsuperscript{3}}
\newunicodechar{₅}{\textsubscript{5}}
\newunicodechar{×}{\ensuremath{\times}}
\newunicodechar{Δ}{\ensuremath{\Delta}}
\newunicodechar{Ω}{\ensuremath{\Omega}}
\newunicodechar{–}{--}
\newunicodechar{—}{---}
\newunicodechar{↑}{\ensuremath{\uparrow}}
\newunicodechar{→}{\ensuremath{\rightarrow}}
\newunicodechar{↓}{\ensuremath{\downarrow}}
\newunicodechar{∝}{\ensuremath{\propto}}
\newunicodechar{‘}{`}
\newunicodechar{’}{'}
\newunicodechar{“}{``}
\newunicodechar{”}{''}

\geometry{margin=2.5cm}

% === 元数据 ===
\title{ 路线 A：构效关系清单（SPR Inventory） }
\date{ 2026-08 }
\author{Pi-Agent（自主文献调研 Agent）}

\begin{document}
\maketitle
\sloppy  % 允许宽松换行，缓解长 URL/DOI 溢出

\section{路线 A：构效关系清单}\label{ux8defux7ebf-aux6784ux6548ux5173ux7cfbux6e05ux5355}

\begin{quote}
\textbf{提交日期}：2026 年 8 月 \textbf{生成脚本}：\texttt{scripts/build\_route\_a\_docs.py} \textbf{覆盖主题}：6 个 \textbf{假设总数}：31 条 \textbf{已验证}：7 条 \textbar{} \textbf{待验证}：21 条 \textbar{} \textbf{待定}：3 条
\end{quote}

\begin{center}\rule{0.5\linewidth}{0.5pt}\end{center}

\subsection{1. 总览}\label{ux603bux89c8}

{\def\LTcaptype{none} % do not increment counter
\begin{longtable}[]{@{}
  >{\raggedright\arraybackslash}p{(\linewidth - 14\tabcolsep) * \real{0.1429}}
  >{\raggedright\arraybackslash}p{(\linewidth - 14\tabcolsep) * \real{0.0857}}
  >{\raggedright\arraybackslash}p{(\linewidth - 14\tabcolsep) * \real{0.1286}}
  >{\raggedright\arraybackslash}p{(\linewidth - 14\tabcolsep) * \real{0.1143}}
  >{\raggedright\arraybackslash}p{(\linewidth - 14\tabcolsep) * \real{0.1143}}
  >{\raggedright\arraybackslash}p{(\linewidth - 14\tabcolsep) * \real{0.1286}}
  >{\raggedright\arraybackslash}p{(\linewidth - 14\tabcolsep) * \real{0.1286}}
  >{\raggedright\arraybackslash}p{(\linewidth - 14\tabcolsep) * \real{0.1571}}@{}}
\toprule\noalign{}
\begin{minipage}[b]{\linewidth}\raggedright
SPR 编号
\end{minipage} & \begin{minipage}[b]{\linewidth}\raggedright
主题
\end{minipage} & \begin{minipage}[b]{\linewidth}\raggedright
构效关系
\end{minipage} & \begin{minipage}[b]{\linewidth}\raggedright
置信度
\end{minipage} & \begin{minipage}[b]{\linewidth}\raggedright
新颖性
\end{minipage} & \begin{minipage}[b]{\linewidth}\raggedright
搜索方法
\end{minipage} & \begin{minipage}[b]{\linewidth}\raggedright
验证状态
\end{minipage} & \begin{minipage}[b]{\linewidth}\raggedright
外部数据库
\end{minipage} \\
\midrule\noalign{}
\endhead
\bottomrule\noalign{}
\endlastfoot
SPR-MOF-01 & MOF CO₂ 捕获 & 双金属MOF-74中金属比例与CO2吸附容量呈倒U型定量关系 & 0.88 & 0.82 & bayesian & pending & materials\_project \\
SPR-MOF-02 & MOF CO₂ 捕获 & 胺功能化MOF在湿态下的CO2捕获容量随相对湿度呈峰值增强 & 0.90 & 0.85 & bayesian & pending & materials\_project \\
SPR-MOF-03 & MOF CO₂ 捕获 & MOF的CO2吸附热（Qst）在25-40 kJ/mol窗口内可实现容量-选择性-再生能耗Pareto最优 & 0.62 & 0.78 & bayesian & pending & materials\_project \\
SPR-MOF-04 & MOF CO₂ 捕获 & MOF-74缺陷类型依赖：缺失配体型缺陷↑OMS/容量 vs 溶剂甲酸盐占据型缺陷↓OMS/容量 & 0.92 & 0.88 & bayesian & validated & materials\_project \\
SPR-MOF-05 & MOF CO₂ 捕获 & MOF在痕量NO2/SO2暴露后CO2吸附容量衰减率与OMS密度和孔道亲水性相关 & 0.63 & 0.90 & bayesian & pending & materials\_project \\
SPR-MOFv4-01 & MOF CO₂ 捕获（e2e v4 重跑） & 双金属 NiCo-MOF-74 组分比例-容量倒U（高斯峰）标度律 & 0.87 & 0.73 & bayesian & validated & materials\_project, oqmd, hmof\_core\_mof, nomad \\
SPR-MOFv4-02 & MOF CO₂ 捕获（e2e v4 重跑） & 湿度-胺 MOF 合作吸附诱导效应标度律（临界RH与胺结构相关） & 0.80 & 0.80 & mcts & pending & --- \\
SPR-MOFv4-03 & MOF CO₂ 捕获（e2e v4 重跑） & 胺结构（碳数/支化/环化）与 CO2/diamine 化学计量上限的标度律 & 0.86 & 0.78 & bayesian & pending & --- \\
SPR-MOFv4-04 & MOF CO₂ 捕获（e2e v4 重跑） & M-MOF-74 金属 d 电子构型-电负性联合描述符预测 CO2 吸附焓 & 0.87 & 0.70 & hybrid & pending & --- \\
SPR-MOFv4-05 & MOF CO₂ 捕获（e2e v4 重跑） & 吸附焓-再生能耗权衡：低 Qst 吸附剂节能但需工艺补偿 & 0.91 & 0.72 & bayesian & validated & materials\_project, oqmd, hmof\_core\_mof, nomad \\
SPR-PVSK-01 & 卤化物钙钛矿 & 带隙-稳定性 trade-off 的定量标度律：窄带隙钙钛矿稳定性系统性下降 & 0.96 & 0.65 & bayesian & validated & materials\_project, oqmd, nomad \\
SPR-PVSK-02 & 卤化物钙钛矿 & 温度-带隙反常标度：dEg/dT\textgreater0 与非谐振动/电声耦合强度相关 & 0.96 & 0.60 & bayesian & pending & --- \\
SPR-PVSK-03 & 卤化物钙钛矿 & 双钙钛矿间接→直接带隙调控：In³⁺/无序/Pb²⁺掺杂的系统效应 & 0.97 & 0.55 & bayesian & validated & materials\_project, oqmd, nomad \\
SPR-PVSK-04 & 卤化物钙钛矿 & 压力-带隙闭合的分段线性标度：MA2PtI6 类体系 dEg/dP 与结构相变耦合 & 0.96 & 0.50 & bayesian & pending & --- \\
SPR-PVSK-05 & 卤化物钙钛矿 & ML 联合预测带隙与稳定性：电离能/分解能描述符的独立性检验 & 0.88 & 0.60 & bayesian & pending & --- \\
SPR-TE-01 & 热电材料 & 纳米晶尺寸调控 Si-Ge-P 超饱和固溶体的晶格热导率与 ZT & 0.87 & 0.80 & bayesian & inconclusive & materials\_project, oqmd, nomad \\
SPR-TE-02 & 热电材料 & 卤素掺杂位点与 n 型 SnSe 载流子浓度及 ZT 的关联 & 0.85 & 0.68 & bayesian & pending & --- \\
SPR-TE-03 & 热电材料 & 最优掺杂浓度随温度上移抑制双极效应的实验验证 & 0.83 & 0.74 & bayesian & pending & --- \\
SPR-TE-04 & 热电材料 & 共振掺杂浓度与 DOS 局部异常及 Seebeck 增强的定量标度 & 0.68 & 0.82 & bayesian & pending & --- \\
SPR-TE-05 & 热电材料 & Yb 填充方钴矿填充分数对晶格热导率与迁移率的定量影响及批次稳定性 & 0.79 & 0.65 & bayesian & inconclusive & materials\_project, oqmd, nomad \\
SPR-NMC-01 & 高镍正极 & 单晶NMC811中Ni氧化态异质性调控位错辅助裂纹萌生 & 0.72 & 0.88 & bayesian & inconclusive & materials\_project, oqmd, nomad \\
SPR-NMC-02 & 高镍正极 & 裂纹电解质润湿通过表面重构与氧释放加速容量衰减 & 0.70 & 0.90 & bayesian & pending & --- \\
SPR-NMC-03 & 高镍正极 & 高镍层状氧化物中Ni含量与容量保持率呈非单调权衡 & 0.68 & 0.82 & bayesian & validated & materials\_project, oqmd, nomad \\
SPR-NMC-04 & 高镍正极 & 涂层材料的Li离子电导率与界面副反应阻抗之间的帕累托最优设计 & 0.65 & 0.78 & bayesian & pending & --- \\
SPR-NMC-05 & 高镍正极 & Ni氧化态升高降低阳离子混排迁移势垒，促进层状结构向岩盐相降解 & 0.63 & 0.85 & bayesian & pending & --- \\
SPR-NMC-06 & 高镍正极 & 水洗引入的质子残留通过近表面Li/H交换加速高镍正极容量衰减与阻抗增长 & 0.65 & 0.85 & bayesian & validated & materials\_project, oqmd, nomad \\
SPR-SE-01 & 固态电解质 & 卤素比例调控卤化物固态电解质室温电导率与湿度稳定性的非线性关系 & 0.72 & 0.70 & bayesian & pending & --- \\
SPR-SE-02 & 固态电解质 & 氧/钼双掺杂对硫银锗矿硫化物电解质电导率与H2S释放量的双目标优化窗口 & 0.79 & 0.75 & bayesian & pending & --- \\
SPR-SE-03 & 固态电解质 & 高熵阳离子无序度与固态电解质室温离子电导率的火山型关系及机器学习预测闭环 & 0.65 & 0.85 & bayesian & pending & --- \\
SPR-SE-04 & 固态电解质 & Lewis酸性基团浓度对聚合物电解质锂离子迁移数与离子电导率权衡的非单调调控 & 0.62 & 0.80 & bayesian & pending & --- \\
SPR-SE-05 & 固态电解质 & 亲锂成核层+电子阻挡层双层界面设计对固态电解质界面电阻的普适降低效应 & 0.66 & 0.75 & bayesian & pending & --- \\
\end{longtable}
}

\subsection{2. 详细清单}\label{ux8be6ux7ec6ux6e05ux5355}

\subsubsection{SPR-MOF-01 --- 双金属MOF-74中金属比例与CO2吸附容量呈倒U型定量关系}\label{spr-mof-01-ux53ccux91d1ux5c5emof-74ux4e2dux91d1ux5c5eux6bd4ux4f8bux4e0eco2ux5438ux9644ux5bb9ux91cfux5448ux5012uux578bux5b9aux91cfux5173ux7cfb}

{\def\LTcaptype{none} % do not increment counter
\begin{longtable}[]{@{}
  >{\raggedright\arraybackslash}p{(\linewidth - 2\tabcolsep) * \real{0.5000}}
  >{\raggedright\arraybackslash}p{(\linewidth - 2\tabcolsep) * \real{0.5000}}@{}}
\toprule\noalign{}
\begin{minipage}[b]{\linewidth}\raggedright
字段
\end{minipage} & \begin{minipage}[b]{\linewidth}\raggedright
内容
\end{minipage} \\
\midrule\noalign{}
\endhead
\bottomrule\noalign{}
\endlastfoot
\textbf{主题/材料体系} & MOF CO₂ 捕获：CoMn-MOF-74、NiCu-MOF-74、FeCu-MOF-74、MIL-101(Cr,Mg) \\
\textbf{性能属性} & CO2吸附容量 \\
\textbf{构效关系陈述} & 金属比例与CO2容量呈倒U型关系，峰值位于特定比例（如1:1或3:1），选择性同步变化；双金属优于任意单金属。 \\
\textbf{置信度} & 0.880 \\
\textbf{新颖性} & 0.82 --- 新知（高新颖性） \\
\textbf{搜索方法} & bayesian \\
\textbf{Best Score} & 0.000 \\
\textbf{验证状态} & pending \\
\textbf{模型比较} & --- \\
\textbf{外部数据库} & materials\_project \\
\textbf{可提取数据点} & 0 点 / 0 种独立材料 \\
\textbf{证据链} & p65, p67, p22, p147, p24 \\
\textbf{已知工作} & 已有文献依据(evidence\_chain 编号: p65, p67, p22, p147)，具体结论需人工/LLM 补写 \\
\textbf{增量贡献} & 相对已确立结论的具体增量待补写(本假设的 expected\_relationship 即拟验证的新规律) \\
\textbf{描述} & 在双金属MOF-74（如CoMn、NiCu、FeCu等）体系中，通过梯度调节两种金属的摩尔比例（0-100\%，步长10\%），系统测量CO2吸附容量和CO2/N2选择性。预期发现容量-比例曲线呈倒U型，且最优比例因金属组合不同而偏移（如CoMn为1:1，NiCu为1:1附近）。该关系源于双金属协同效应：异金属节点改变电子密度和不饱和配位环境，从而优化CO2结合能，但过量第二金属可能导致结构扭曲或相分离，削弱性能。 \\
\end{longtable}
}

{[}数值验证{]} 以下数值未在文献中查证: 10\% \textbar{}

\subsubsection{SPR-MOF-02 --- 胺功能化MOF在湿态下的CO2捕获容量随相对湿度呈峰值增强}\label{spr-mof-02-ux80faux529fux80fdux5316mofux5728ux6e7fux6001ux4e0bux7684co2ux6355ux83b7ux5bb9ux91cfux968fux76f8ux5bf9ux6e7fux5ea6ux5448ux5cf0ux503cux589eux5f3a}

{\def\LTcaptype{none} % do not increment counter
\begin{longtable}[]{@{}
  >{\raggedright\arraybackslash}p{(\linewidth - 2\tabcolsep) * \real{0.5000}}
  >{\raggedright\arraybackslash}p{(\linewidth - 2\tabcolsep) * \real{0.5000}}@{}}
\toprule\noalign{}
\begin{minipage}[b]{\linewidth}\raggedright
字段
\end{minipage} & \begin{minipage}[b]{\linewidth}\raggedright
内容
\end{minipage} \\
\midrule\noalign{}
\endhead
\bottomrule\noalign{}
\endlastfoot
\textbf{主题/材料体系} & MOF CO₂ 捕获：胺接枝Mg-MOF-74、MOF@MOF核壳材料、疏水功能化胺浸渍树脂、Mg-MOF-74 \\
\textbf{性能属性} & 湿态CO2吸附容量 \\
\textbf{构效关系陈述} & 胺功能化材料的CO2容量随RH先增后减，峰值在中等RH区域；OMS型材料的容量随RH增加而单调衰减。 \\
\textbf{置信度} & 0.900 \\
\textbf{新颖性} & 0.85 --- 新知（高新颖性） \\
\textbf{搜索方法} & bayesian \\
\textbf{Best Score} & 0.000 \\
\textbf{验证状态} & pending \\
\textbf{模型比较} & --- \\
\textbf{外部数据库} & materials\_project \\
\textbf{可提取数据点} & 0 点 / 0 种独立材料 \\
\textbf{证据链} & p210, p224, p201, p186, p182 \\
\textbf{已知工作} & 已有文献依据(evidence\_chain 编号: p210, p224, p201, p186)，具体结论需人工/LLM 补写 \\
\textbf{增量贡献} & 相对已确立结论的具体增量待补写(本假设的 expected\_relationship 即拟验证的新规律) \\
\textbf{描述} & 对比OMS型MOF（如Mg-MOF-74）和胺功能化MOF（如胺接枝MOF或MOF@MOF核壳结构）在0-90\% RH范围内的CO2吸附容量。预期胺型材料在中等湿度（如40-60\% RH）出现容量峰值，因水分子促进氨基甲酸酯形成并缓解传质限制；而OMS型材料随RH增加容量单调下降，因水优先占据开放金属位点。该假设可统一现有文献中水促进/抑制的矛盾结论。 \\
\end{longtable}
}

\subsubsection{SPR-MOF-03 --- MOF的CO2吸附热（Qst）在25-40 kJ/mol窗口内可实现容量-选择性-再生能耗Pareto最优}\label{spr-mof-03-mofux7684co2ux5438ux9644ux70edqstux572825-40-kjmolux7a97ux53e3ux5185ux53efux5b9eux73b0ux5bb9ux91cf-ux9009ux62e9ux6027-ux518dux751fux80fdux8017paretoux6700ux4f18}

{\def\LTcaptype{none} % do not increment counter
\begin{longtable}[]{@{}
  >{\raggedright\arraybackslash}p{(\linewidth - 2\tabcolsep) * \real{0.5000}}
  >{\raggedright\arraybackslash}p{(\linewidth - 2\tabcolsep) * \real{0.5000}}@{}}
\toprule\noalign{}
\begin{minipage}[b]{\linewidth}\raggedright
字段
\end{minipage} & \begin{minipage}[b]{\linewidth}\raggedright
内容
\end{minipage} \\
\midrule\noalign{}
\endhead
\bottomrule\noalign{}
\endlastfoot
\textbf{主题/材料体系} & MOF CO₂ 捕获：MOF-74(Ni)、MOF-74(Mg)、胺化SBA-15、锂掺杂MIL-101、UiO-66 \\
\textbf{性能属性} & CO2吸附热（Qst） \\
\textbf{构效关系陈述} & Qst与选择性呈正相关，与再生能耗也正相关；存在Qst甜点区（25-40 kJ/mol），在此区间容量-选择性-再生能耗三者可达Pareto最优。 \\
\textbf{置信度} & 0.620 \\
\textbf{新颖性} & 0.78 --- 新知（中等新颖性） \\
\textbf{搜索方法} & bayesian \\
\textbf{Best Score} & 0.000 \\
\textbf{验证状态} & pending \\
\textbf{模型比较} & --- \\
\textbf{外部数据库} & materials\_project \\
\textbf{可提取数据点} & 0 点 / 0 种独立材料 \\
\textbf{证据链} & p27, p117, p129, p116 \\
\textbf{已知工作} & 已有文献依据(evidence\_chain 编号: p27, p117, p129, p116)，具体结论需人工/LLM 补写 \\
\textbf{增量贡献} & 相对已确立结论的具体增量待补写(本假设的 expected\_relationship 即拟验证的新规律) \\
\textbf{描述} & 选取50种以上具有不同Qst值的MOF材料（涵盖MOF-74系列、UiO系列、ZIF系列等），测量其CO2容量、CO2/N2选择性和再生温度/能耗。预期Qst在29-40 kJ/mol（尤其约29 kJ/mol）时，材料同时具备较高容量和选择性，且再生能耗适中；过高的Qst（\textgreater50）虽增强容量但再生困难，过低Qst（\textless20）选择性不足。该假设将建立容量-选择性-再生能三维Pareto前沿，为吸附剂筛选提供热力学约束。 \\
\end{longtable}
}

{[}数值验证{]} 以下数值未在文献中查证: 25-40 kJ/mol, 29-40 kJ/mol, 29 kJ/mol, 25-40 kJ/mol \textbar{}

\subsubsection{SPR-MOF-04 --- MOF-74缺陷类型依赖：缺失配体型缺陷↑OMS/容量 vs 溶剂甲酸盐占据型缺陷↓OMS/容量}\label{spr-mof-04-mof-74ux7f3aux9677ux7c7bux578bux4f9dux8d56ux7f3aux5931ux914dux4f53ux578bux7f3aux9677omsux5bb9ux91cf-vs-ux6eb6ux5242ux7532ux9178ux76d0ux5360ux636eux578bux7f3aux9677omsux5bb9ux91cf}

{\def\LTcaptype{none} % do not increment counter
\begin{longtable}[]{@{}
  >{\raggedright\arraybackslash}p{(\linewidth - 2\tabcolsep) * \real{0.5000}}
  >{\raggedright\arraybackslash}p{(\linewidth - 2\tabcolsep) * \real{0.5000}}@{}}
\toprule\noalign{}
\begin{minipage}[b]{\linewidth}\raggedright
字段
\end{minipage} & \begin{minipage}[b]{\linewidth}\raggedright
内容
\end{minipage} \\
\midrule\noalign{}
\endhead
\bottomrule\noalign{}
\endlastfoot
\textbf{主题/材料体系} & MOF CO₂ 捕获：缺陷MOF-74(Co)、缺陷MOF-74(Mg)、缺陷MOF-74(Ni) \\
\textbf{性能属性} & 开放金属位点密度及CO2吸附容量 \\
\textbf{构效关系陈述} & 缺陷类型依赖（非普适线性）：缺失配体型→OMS/容量正相关（先升后降拐点）；溶剂甲酸盐占据型→OMS/容量负相关。忽略缺陷类型将导致方向性误判。 \\
\textbf{置信度} & 0.920 \\
\textbf{新颖性} & 0.88 --- 新知（高新颖性） \\
\textbf{搜索方法} & bayesian \\
\textbf{Best Score} & 0.000 \\
\textbf{验证状态} & validated \\
\textbf{模型比较} & --- \\
\textbf{外部数据库} & materials\_project \\
\textbf{可提取数据点} & 0 点 / 0 种独立材料 \\
\textbf{证据链} & p13, p7, p192, p9, r10s1\_a9ed68d1734e, r10s0\_1fbbe940f717, r10s1\_821919f491c8, r10s2\_0341d732c860 \\
\textbf{已知工作} & 已有文献依据(evidence\_chain 编号: p13, p7, p192, p9)，具体结论需人工/LLM 补写 \\
\textbf{增量贡献} & 相对已确立结论的具体增量待补写(本假设的 expected\_relationship 即拟验证的新规律) \\
\textbf{描述} & 区分两类缺陷通道对CO2吸附的方向性影响（九轮缺陷类型二分，Nature Commun 2023）：(1) 缺失配体型（HBDC比例0-50\%）------配体缺失暴露额外金属位点，OMS密度↑、CO2容量↑（但超过临界比例后结构稳定性下降、容量回落，临界值待实验确定）；(2) 溶剂衍生甲酸盐型（DMF分解甲酸配位金属）------甲酸盐额外M-O键部分消除OMS，N2/CO2吸附随缺陷浓度定量下降（r10s1\_a9ed68d1734e）。十轮新增旁证：水桥可替代缺陷增强胺功能化UiO-66的CO2吸附（r10s0\_1fbbe940f717）、Ru-MOF不同缺陷类型对吸附方向性影响不同（r10s1\_8219 \\
\end{longtable}
}

\subsubsection{SPR-MOF-05 --- MOF在痕量NO2/SO2暴露后CO2吸附容量衰减率与OMS密度和孔道亲水性相关}\label{spr-mof-05-mofux5728ux75d5ux91cfno2so2ux66b4ux9732ux540eco2ux5438ux9644ux5bb9ux91cfux8870ux51cfux7387ux4e0eomsux5bc6ux5ea6ux548cux5b54ux9053ux4eb2ux6c34ux6027ux76f8ux5173}

{\def\LTcaptype{none} % do not increment counter
\begin{longtable}[]{@{}
  >{\raggedright\arraybackslash}p{(\linewidth - 2\tabcolsep) * \real{0.5000}}
  >{\raggedright\arraybackslash}p{(\linewidth - 2\tabcolsep) * \real{0.5000}}@{}}
\toprule\noalign{}
\begin{minipage}[b]{\linewidth}\raggedright
字段
\end{minipage} & \begin{minipage}[b]{\linewidth}\raggedright
内容
\end{minipage} \\
\midrule\noalign{}
\endhead
\bottomrule\noalign{}
\endlastfoot
\textbf{主题/材料体系} & MOF CO₂ 捕获：Mg-MOF-74、HKUST-1、ZIF-8、MIL-101(Cr) \\
\textbf{性能属性} & CO2吸附容量保持率 \\
\textbf{构效关系陈述} & 暴露NO2/SO2后，CO2容量保持率与OMS密度呈负相关，与孔道疏水性呈正相关；ZIF-8保持率最高，MOF-74最低。 \\
\textbf{置信度} & 0.630 \\
\textbf{新颖性} & 0.90 --- 新知（高新颖性） \\
\textbf{搜索方法} & bayesian \\
\textbf{Best Score} & 0.000 \\
\textbf{验证状态} & pending \\
\textbf{模型比较} & --- \\
\textbf{外部数据库} & materials\_project \\
\textbf{可提取数据点} & 0 点 / 0 种独立材料 \\
\textbf{证据链} & p216, p52, p55 \\
\textbf{已知工作} & 已有文献依据(evidence\_chain 编号: p216, p52, p55)，具体结论需人工/LLM 补写 \\
\textbf{增量贡献} & 相对已确立结论的具体增量待补写(本假设的 expected\_relationship 即拟验证的新规律) \\
\textbf{描述} & 选取MOF-74(Mg)、ZIF-8、HKUST-1、MIL-101(Cr)等代表性MOF，在含1000 ppm NO2或SO2的模拟烟气/空气中暴露不同时间（0-24h），测量CO2吸附容量变化及结构变化。预期含高密度OMS的MOF（如MOF-74、HKUST-1）衰减更快，因杂质气体与金属位点强结合或反应；而ZIF-8因疏水性和配位饱和更稳定。该假设为MOF在真实含杂质环境中的适用性提供筛选依据。 \\
\end{longtable}
}

{[}数值验证{]} 以下数值未在文献中查证: 1000 ppm, 0-24h \textbar{}

\subsubsection{SPR-MOFv4-01 --- 双金属 NiCo-MOF-74 组分比例-容量倒U（高斯峰）标度律}\label{spr-mofv4-01-ux53ccux91d1ux5c5e-nico-mof-74-ux7ec4ux5206ux6bd4ux4f8b-ux5bb9ux91cfux5012uux9ad8ux65afux5cf0ux6807ux5ea6ux5f8b}

{\def\LTcaptype{none} % do not increment counter
\begin{longtable}[]{@{}
  >{\raggedright\arraybackslash}p{(\linewidth - 2\tabcolsep) * \real{0.5000}}
  >{\raggedright\arraybackslash}p{(\linewidth - 2\tabcolsep) * \real{0.5000}}@{}}
\toprule\noalign{}
\begin{minipage}[b]{\linewidth}\raggedright
字段
\end{minipage} & \begin{minipage}[b]{\linewidth}\raggedright
内容
\end{minipage} \\
\midrule\noalign{}
\endhead
\bottomrule\noalign{}
\endlastfoot
\textbf{主题/材料体系} & MOF CO₂ 捕获（e2e v4 重跑）：NiCo-MOF-74、Co-MOF-74、Ni-MOF-74、Ni1Co6-MOF-74、Ni6Co1-MOF-74 \\
\textbf{性能属性} & CO2 容量 (mmol/g) @ 0C 1bar \\
\textbf{构效关系陈述} & C(x) = A*exp(-(x-mu)\textsuperscript{2/(2*sigma}2)) + B, 峰值 x\_mu \textasciitilde{} 0.5, 容量增益 \textasciitilde5.6 mmol/g \\
\textbf{置信度} & 0.868 \\
\textbf{新颖性} & 0.73 --- 新知（中等新颖性） \\
\textbf{搜索方法} & bayesian \\
\textbf{Best Score} & 0.000 \\
\textbf{验证状态} & validated \\
\textbf{模型比较} & \{`verdict': `insufficient', `reason': `候选或经典 R² 缺失，无法判定', `delta\_r2': None, `f\_supported': False, `ci\_supported': False\} \\
\textbf{外部数据库} & materials\_project, oqmd, hmof\_core\_mof, nomad \\
\textbf{可提取数据点} & 0 点 / 0 种独立材料 \\
\textbf{证据链} & （证据链条目存在，但需清理格式） \\
\textbf{已知工作} & 单金属 MOF-74 Qst 排序（Caskey 2008）已确立；但双金属组分-容量连续标度函数未见报道 \\
\textbf{增量贡献} & 给出倒U的定量函数形式（高斯参数 mu/sigma/A/B）并预测峰值位置与曲率随金属对的迁移规律 \\
\textbf{描述} & NiCo-MOF-74 系列 CO2 容量随 Ni/(Ni+Co) 比例呈倒U/高斯型：C(x)=A*exp(-(x-mu)\textsuperscript{2/(2*sigma}2))+B。5 点实验数据（Co 5.03 / Ni1Co6 6.40 / Ni1Co1 8.30 / Ni6Co1 3.62 / Ni 3.99 mmol/g @0C 1bar）显示峰值在 x\textasciitilde0.5。假设：倒U曲率与峰值位置由金属对 d-d 电子构型差异决定，可用于预测其他 d-d 金属对（Ni/Mn、Co/Mn）的最佳组分。 \\
\end{longtable}
}

\subsubsection{SPR-MOFv4-02 --- 湿度-胺 MOF 合作吸附诱导效应标度律（临界RH与胺结构相关）}\label{spr-mofv4-02-ux6e7fux5ea6-ux80fa-mof-ux5408ux4f5cux5438ux9644ux8bf1ux5bfcux6548ux5e94ux6807ux5ea6ux5f8bux4e34ux754crhux4e0eux80faux7ed3ux6784ux76f8ux5173}

{\def\LTcaptype{none} % do not increment counter
\begin{longtable}[]{@{}
  >{\raggedright\arraybackslash}p{(\linewidth - 2\tabcolsep) * \real{0.5000}}
  >{\raggedright\arraybackslash}p{(\linewidth - 2\tabcolsep) * \real{0.5000}}@{}}
\toprule\noalign{}
\begin{minipage}[b]{\linewidth}\raggedright
字段
\end{minipage} & \begin{minipage}[b]{\linewidth}\raggedright
内容
\end{minipage} \\
\midrule\noalign{}
\endhead
\bottomrule\noalign{}
\endlastfoot
\textbf{主题/材料体系} & MOF CO₂ 捕获（e2e v4 重跑）：mmen-Mg2(dobpdc)、dmpn-Mg2(dobpdc)、e-2-Mg2(dobpdc)、een-MOF、IRMOF-74-III \\
\textbf{性能属性} & 诱导效应强度 I (0-1), CO2 吸附速率 (相对) \\
\textbf{构效关系陈述} & I(RH) = I0 - k\emph{RH (线性) 或 I(RH) = I0}exp(-RH/tau)；临界 RHc = I0/k 随胺碳数↑而↑ \\
\textbf{置信度} & 0.798 \\
\textbf{新颖性} & 0.80 --- 新知（高新颖性） \\
\textbf{搜索方法} & mcts \\
\textbf{Best Score} & 0.000 \\
\textbf{验证状态} & pending \\
\textbf{模型比较} & --- \\
\textbf{外部数据库} & --- \\
\textbf{可提取数据点} & 0 点 / 0 种独立材料 \\
\textbf{证据链} & （证据链条目存在，但需清理格式） \\
\textbf{已知工作} & Marshall 2024 定性确立 RH 效应；跨材料定量标度未见报道 \\
\textbf{增量贡献} & 将诱导效应定量化为 I(RH) 函数并给出临界 RHc 与胺结构/孔径的关联 \\
\textbf{描述} & diamine-appended MOF（mmen-Mg2(dobpdc) 等）在低 RH 下出现诱导效应（合作链），高 RH 下诱导效应消失、吸附速率增加（LKT 理论）。假设：诱导效应强度 I(RH) 随 RH 单调下降（可拟合成线性或指数衰减），临界 RHc（I=0 的截距）随胺链长增加而升高、随骨架孔径增加而降低。该标度律可用于预判 DAC/湿烟气工况下胺 MOF 的性能。 \\
\end{longtable}
}

\subsubsection{SPR-MOFv4-03 --- 胺结构（碳数/支化/环化）与 CO2/diamine 化学计量上限的标度律}\label{spr-mofv4-03-ux80faux7ed3ux6784ux78b3ux6570ux652fux5316ux73afux5316ux4e0e-co2diamine-ux5316ux5b66ux8ba1ux91cfux4e0aux9650ux7684ux6807ux5ea6ux5f8b}

{\def\LTcaptype{none} % do not increment counter
\begin{longtable}[]{@{}
  >{\raggedright\arraybackslash}p{(\linewidth - 2\tabcolsep) * \real{0.5000}}
  >{\raggedright\arraybackslash}p{(\linewidth - 2\tabcolsep) * \real{0.5000}}@{}}
\toprule\noalign{}
\begin{minipage}[b]{\linewidth}\raggedright
字段
\end{minipage} & \begin{minipage}[b]{\linewidth}\raggedright
内容
\end{minipage} \\
\midrule\noalign{}
\endhead
\bottomrule\noalign{}
\endlastfoot
\textbf{主题/材料体系} & MOF CO₂ 捕获（e2e v4 重跑）：en-Mg2(dobpdc)、mmen-Mg2(dobpdc)、e-2-Mg2(dobpdc)、dmpn-Mg2(dobpdc)、pip2-Mg2(dobpdc) \\
\textbf{性能属性} & CO2/diamine 化学计量 (mol/mol) \\
\textbf{构效关系陈述} & 化学计量 S = 1.0 + f(环化度)，环状双胺（pip2）可达 \textasciitilde1.5 \\
\textbf{置信度} & 0.855 \\
\textbf{新颖性} & 0.78 --- 新知（中等新颖性） \\
\textbf{搜索方法} & bayesian \\
\textbf{Best Score} & 0.000 \\
\textbf{验证状态} & pending \\
\textbf{模型比较} & --- \\
\textbf{外部数据库} & --- \\
\textbf{可提取数据点} & 0 点 / 0 种独立材料 \\
\textbf{证据链} & （证据链条目存在，但需清理格式） \\
\textbf{已知工作} & 单一胺化学吸附机理已确立（Forse 2018）；双机理突破上限仅 pip2 单点报道 \\
\textbf{增量贡献} & 建立胺拓扑结构与化学计量上限的定量关系，指导高容量胺设计 \\
\textbf{描述} & diamine-appended Mg2(dobpdc) 化学吸附化学计量传统上限为 1.0 CO2/diamine（铵盐-氨基甲酸酯链），但 pip2（1-(2-氨基乙基)哌啶，双环胺）实现 \textasciitilde1.5 CO2/diamine 两步吸附。假设：胺的支化/环化程度（拓扑复杂度）与化学计量上限正相关------含环状仲胺的体系可通过物理+化学双机理突破化学计量上限。 \\
\end{longtable}
}

\subsubsection{SPR-MOFv4-04 --- M-MOF-74 金属 d 电子构型-电负性联合描述符预测 CO2 吸附焓}\label{spr-mofv4-04-m-mof-74-ux91d1ux5c5e-d-ux7535ux5b50ux6784ux578b-ux7535ux8d1fux6027ux8054ux5408ux63cfux8ff0ux7b26ux9884ux6d4b-co2-ux5438ux9644ux7113}

{\def\LTcaptype{none} % do not increment counter
\begin{longtable}[]{@{}
  >{\raggedright\arraybackslash}p{(\linewidth - 2\tabcolsep) * \real{0.5000}}
  >{\raggedright\arraybackslash}p{(\linewidth - 2\tabcolsep) * \real{0.5000}}@{}}
\toprule\noalign{}
\begin{minipage}[b]{\linewidth}\raggedright
字段
\end{minipage} & \begin{minipage}[b]{\linewidth}\raggedright
内容
\end{minipage} \\
\midrule\noalign{}
\endhead
\bottomrule\noalign{}
\endlastfoot
\textbf{主题/材料体系} & MOF CO₂ 捕获（e2e v4 重跑）：Mg-MOF-74、Fe-MOF-74、Co-MOF-74、Ni-MOF-74、Zn-MOF-74、Ca-MOF-74、Mn-MOF-74、Cu-MOF-74、Ti-MOF-74 \\
\textbf{性能属性} & Qst\_CO2 (kJ/mol) \\
\textbf{构效关系陈述} & Qst = a\emph{Nd + b}chi + c*chi\^{}2 + d；Mg(0d, 1.31)-\textgreater47, Fe(6d,1.83)-\textgreater36, Co(7d,1.88)-\textgreater37, Ni(8d,1.91)-\textgreater41, Zn(10d,1.65)-\textgreater29 \\
\textbf{置信度} & 0.866 \\
\textbf{新颖性} & 0.70 --- 新知（中等新颖性） \\
\textbf{搜索方法} & hybrid \\
\textbf{Best Score} & 0.000 \\
\textbf{验证状态} & pending \\
\textbf{模型比较} & \{`verdict': `insufficient', `reason': `候选或经典 R² 缺失，无法判定', `delta\_r2': None, `f\_supported': False, `ci\_supported': False\} \\
\textbf{外部数据库} & --- \\
\textbf{可提取数据点} & 0 点 / 0 种独立材料 \\
\textbf{证据链} & （证据链条目存在，但需清理格式） \\
\textbf{已知工作} & Caskey/Koh 确立金属种类排序；双描述符连续函数形式未见报道 \\
\textbf{增量贡献} & 给出 Qst=f(Nd, chi) 的定量函数并外推预测 Ca/Mn/Cu/Ti \\
\textbf{描述} & M-MOF-74 系列 CO2 吸附焓 Qst 由金属阳离子电子构型决定：d 电子数与电负性联合描述符可预测 Qst（v3 拟合 R²=0.9855）。假设：Qst = a\emph{Nd + b}chi + c*chi\^{}2 + d（二次多项式），该标度可外推预测未实验金属（如 Ca、Mn、Cu、Ti）的 Qst，指导金属取代筛选。 \\
\end{longtable}
}

\subsubsection{SPR-MOFv4-05 --- 吸附焓-再生能耗权衡：低 Qst 吸附剂节能但需工艺补偿}\label{spr-mofv4-05-ux5438ux9644ux7113-ux518dux751fux80fdux8017ux6743ux8861ux4f4e-qst-ux5438ux9644ux5242ux8282ux80fdux4f46ux9700ux5de5ux827aux8865ux507f}

{\def\LTcaptype{none} % do not increment counter
\begin{longtable}[]{@{}
  >{\raggedright\arraybackslash}p{(\linewidth - 2\tabcolsep) * \real{0.5000}}
  >{\raggedright\arraybackslash}p{(\linewidth - 2\tabcolsep) * \real{0.5000}}@{}}
\toprule\noalign{}
\begin{minipage}[b]{\linewidth}\raggedright
字段
\end{minipage} & \begin{minipage}[b]{\linewidth}\raggedright
内容
\end{minipage} \\
\midrule\noalign{}
\endhead
\bottomrule\noalign{}
\endlastfoot
\textbf{主题/材料体系} & MOF CO₂ 捕获（e2e v4 重跑）：Mg-MOF-74、ED@MOF-520、mmen-Mg2(dobpdc)、CALF-20、MOF纳米复合(磁感应) \\
\textbf{性能属性} & 再生能耗 E\_reg (MJ/kg CO2), Qst (kJ/mol) \\
\textbf{构效关系陈述} & E\_reg ≈ alpha*Qst + beta；磁感应再生 1.29 MJ/kg 为工艺侧节能旁证 \\
\textbf{置信度} & 0.907 \\
\textbf{新颖性} & 0.72 --- 新知（中等新颖性） \\
\textbf{搜索方法} & bayesian \\
\textbf{Best Score} & 0.000 \\
\textbf{验证状态} & validated \\
\textbf{模型比较} & --- \\
\textbf{外部数据库} & materials\_project, oqmd, hmof\_core\_mof, nomad \\
\textbf{可提取数据点} & 0 点 / 0 种独立材料 \\
\textbf{证据链} & （证据链条目存在，但需清理格式） \\
\textbf{已知工作} & 单材料能耗报道多；跨材料 Qst-E\_reg 线性标度未见 \\
\textbf{增量贡献} & 建立 Qst-E\_reg 跨材料标度并给出最优 Qst 窗口 \\
\textbf{描述} & MOF CO2 捕获再生能耗与吸附焓强相关：高 Qst（如 Mg-MOF-74 47 kJ/mol）绑定强但再生能耗高；低 Qst（如 ED@MOF-520 29 kJ/mol）再生节能但需低温/低压工况。假设：再生能耗 E\_reg 与 Qst 近似线性正相关（E\_reg ≈ alpha*Qst + beta），且存在最优 Qst 区间（\textasciitilde30-40 kJ/mol）使总能耗最低------可指导吸附剂与工艺的联合设计。 \\
\end{longtable}
}

\subsubsection{SPR-PVSK-01 --- 带隙-稳定性 trade-off 的定量标度律：窄带隙钙钛矿稳定性系统性下降}\label{spr-pvsk-01-ux5e26ux9699-ux7a33ux5b9aux6027-trade-off-ux7684ux5b9aux91cfux6807ux5ea6ux5f8bux7a84ux5e26ux9699ux9499ux949bux77ffux7a33ux5b9aux6027ux7cfbux7edfux6027ux4e0bux964d}

{\def\LTcaptype{none} % do not increment counter
\begin{longtable}[]{@{}
  >{\raggedright\arraybackslash}p{(\linewidth - 2\tabcolsep) * \real{0.5000}}
  >{\raggedright\arraybackslash}p{(\linewidth - 2\tabcolsep) * \real{0.5000}}@{}}
\toprule\noalign{}
\begin{minipage}[b]{\linewidth}\raggedright
字段
\end{minipage} & \begin{minipage}[b]{\linewidth}\raggedright
内容
\end{minipage} \\
\midrule\noalign{}
\endhead
\bottomrule\noalign{}
\endlastfoot
\textbf{主题/材料体系} & 卤化物钙钛矿：FASnI3、MAPbI3、CsSnI3、Cs2AgBiBr6、AgInI4、CsPbBr3、CH3NH3BaI3、K2SnGeI6 \\
\textbf{性能属性} & band gap \\
\textbf{构效关系陈述} & 带隙 Eg 与稳定性呈正相关（窄带隙不稳定），符合幂律/指数型 trade-off；Eg \textless{} 1.6 eV 体系稳定性骤降 \\
\textbf{置信度} & 0.964 \\
\textbf{新颖性} & 0.65 --- 新知（中等新颖性） \\
\textbf{搜索方法} & bayesian \\
\textbf{Best Score} & 0.000 \\
\textbf{验证状态} & validated \\
\textbf{模型比较} & \{`verdict': `no\_improvement', `reason': `候选 R²=0.0214 与经典 R²=-0.0000 差距不足（ΔR²=+0.0214，\textless0.05 阈值），无显著提升', `delta\_r2': 0.0214, `f\_supported': False, `ci\_supported': False\} \\
\textbf{外部数据库} & materials\_project, oqmd, nomad \\
\textbf{可提取数据点} & 8 点 / 8 种独立材料 \\
\textbf{证据链} & （证据链条目存在，但需清理格式） \\
\textbf{已知工作} & Cu-In 钙钛矿论文（jacs.7b02120）提出稳定性与光电性能可能互相矛盾；阳离子变换（jacs.6b09645）设计稳定无铅体系 \\
\textbf{增量贡献} & 将 trade-off 从定性矛盾提升为定量标度律（Eg-分解能幂律/指数关系），并给出 Pareto 前沿边界 \\
\textbf{描述} & 基于表 D 的 8 材料跨体系数据（FASnI3 1.40 eV / MAPbI3 1.55 eV / CsSnI3 1.30 eV / Cs2AgBiBr6 1.72 eV / AgInI4 1.72 eV / CsPbBr3 2.30 eV / CH3NH3BaI3 3.87 eV / K2SnGeI6 0.64 eV），检验带隙 Eg 与稳定性（分解能/氧化稳定性）之间是否存在单调 trade-off 标度律。预期窄带隙体系（Sn²⁺、低价态）因氧化/分解倾向而系统性不稳定，稳定体系带隙偏高------即存在类似 Pareto 前沿的 Eg-稳定性联合边界。 \\
\end{longtable}
}

\subsubsection{SPR-PVSK-02 --- 温度-带隙反常标度：dEg/dT\textgreater0 与非谐振动/电声耦合强度相关}\label{spr-pvsk-02-ux6e29ux5ea6-ux5e26ux9699ux53cdux5e38ux6807ux5ea6degdt0-ux4e0eux975eux8c10ux632fux52a8ux7535ux58f0ux8026ux5408ux5f3aux5ea6ux76f8ux5173}

{\def\LTcaptype{none} % do not increment counter
\begin{longtable}[]{@{}
  >{\raggedright\arraybackslash}p{(\linewidth - 2\tabcolsep) * \real{0.5000}}
  >{\raggedright\arraybackslash}p{(\linewidth - 2\tabcolsep) * \real{0.5000}}@{}}
\toprule\noalign{}
\begin{minipage}[b]{\linewidth}\raggedright
字段
\end{minipage} & \begin{minipage}[b]{\linewidth}\raggedright
内容
\end{minipage} \\
\midrule\noalign{}
\endhead
\bottomrule\noalign{}
\endlastfoot
\textbf{主题/材料体系} & 卤化物钙钛矿：CsSnI3、CsPbBr3、MAPbI3 \\
\textbf{性能属性} & band gap \\
\textbf{构效关系陈述} & 带隙重正化 ΔEg 随温度 T 按非谐机制增长（ΔEg ∝ T\^{}α 或指数），Slack/Varshni 经典模型因忽略非谐项而失效 \\
\textbf{置信度} & 0.964 \\
\textbf{新颖性} & 0.60 --- 增量新知 \\
\textbf{搜索方法} & bayesian \\
\textbf{Best Score} & 0.000 \\
\textbf{验证状态} & pending \\
\textbf{模型比较} & \{`verdict': `no\_improvement', `reason': `候选 R²=0.0214 与经典 R²=-0.0000 差距不足（ΔR²=+0.0214，\textless0.05 阈值），无显著提升', `delta\_r2': 0.0214, `f\_supported': False, `ci\_supported': False\} \\
\textbf{外部数据库} & --- \\
\textbf{可提取数据点} & 4 点 / 3 种独立材料 \\
\textbf{证据链} & （证据链条目存在，但需清理格式） \\
\textbf{已知工作} & Slack 1971 带隙-温度模型、Varshni 方程广泛应用于半导体；Patrick \& Thygesen (2015) 已发现 CsSnI3 非谐稳定化 \\
\textbf{增量贡献} & 用文献数值系统对比 Slack 模型与候选模型在卤化物钙钛矿上的 R²/RMSE，定量证明经典模型失效并给出非谐标度形式 \\
\textbf{描述} & 铅卤化物钙钛矿带隙随温度降低而降低（反常 dEg/dT\textgreater0，p39），CsSnI3 振动重正化打开带隙 0.11 eV@300K / 0.24 eV@500K（p36），CsPbBr3 非谐贡献 450 meV@425K（p38）。假设 dEg/dT 的符号与幅度由非谐声子/电声耦合（rattler 模式、八面体倾斜）主导，而非经典 Varshni/Slack 模型的热膨胀项。用表 A 数据检验 Slack 模型（经典）与候选模型（幂律/指数）的拟合优劣。 \\
\end{longtable}
}

\subsubsection{SPR-PVSK-03 --- 双钙钛矿间接→直接带隙调控：In³⁺/无序/Pb²⁺掺杂的系统效应}\label{spr-pvsk-03-ux53ccux9499ux949bux77ffux95f4ux63a5ux76f4ux63a5ux5e26ux9699ux8c03ux63a7inuxb3ux65e0ux5e8fpbuxb2ux63baux6742ux7684ux7cfbux7edfux6548ux5e94}

{\def\LTcaptype{none} % do not increment counter
\begin{longtable}[]{@{}
  >{\raggedright\arraybackslash}p{(\linewidth - 2\tabcolsep) * \real{0.5000}}
  >{\raggedright\arraybackslash}p{(\linewidth - 2\tabcolsep) * \real{0.5000}}@{}}
\toprule\noalign{}
\begin{minipage}[b]{\linewidth}\raggedright
字段
\end{minipage} & \begin{minipage}[b]{\linewidth}\raggedright
内容
\end{minipage} \\
\midrule\noalign{}
\endhead
\bottomrule\noalign{}
\endlastfoot
\textbf{主题/材料体系} & 卤化物钙钛矿：Cs2AgBiBr6、Cs2InAgCl6、AgInI4 \\
\textbf{性能属性} & band gap \\
\textbf{构效关系陈述} & In/无序度/掺杂浓度 x 增加 → 带隙类型由间接转直接，带隙大小单调变化（存在转变阈值 x\_c） \\
\textbf{置信度} & 0.966 \\
\textbf{新颖性} & 0.55 --- 增量新知 \\
\textbf{搜索方法} & bayesian \\
\textbf{Best Score} & 0.000 \\
\textbf{验证状态} & validated \\
\textbf{模型比较} & \{`verdict': `insufficient', `reason': `候选或经典 R² 缺失，无法判定', `delta\_r2': None, `f\_supported': False, `ci\_supported': False\} \\
\textbf{外部数据库} & materials\_project, oqmd, nomad \\
\textbf{可提取数据点} & 4 点 / 3 种独立材料 \\
\textbf{证据链} & （证据链条目存在，但需清理格式） \\
\textbf{已知工作} & PhysRevMaterials.2.055401 (p11) 报道 Pb²⁺ 掺杂间接→直接转变；1611.05426v2 (p1) 预测 Cs2InAgCl6 直接带隙 \\
\textbf{增量贡献} & 统一比较 In 替代/无序/Pb 掺杂三条调控路径，提出带隙类型转变的组成阈值规律 \\
\textbf{描述} & Cs2AgBiBr6 为间接带隙 1.72-1.98 eV（p14/p17）；Pb²⁺ 微量掺杂可致间接→直接转变（p11）；Cs2InAgCl6 理论预测直接带隙（p1）；AgInI4 直接带隙 1.72 eV（p70）。假设在双钙钛矿母体中，In 替代/无序增强/掺杂可系统性地将间接带隙转为直接带隙并降低带隙大小，且存在组成-带隙类型转变的临界阈值。 \\
\end{longtable}
}

\subsubsection{SPR-PVSK-04 --- 压力-带隙闭合的分段线性标度：MA2PtI6 类体系 dEg/dP 与结构相变耦合}\label{spr-pvsk-04-ux538bux529b-ux5e26ux9699ux95edux5408ux7684ux5206ux6bb5ux7ebfux6027ux6807ux5ea6ma2pti6-ux7c7bux4f53ux7cfb-degdp-ux4e0eux7ed3ux6784ux76f8ux53d8ux8026ux5408}

{\def\LTcaptype{none} % do not increment counter
\begin{longtable}[]{@{}
  >{\raggedright\arraybackslash}p{(\linewidth - 2\tabcolsep) * \real{0.5000}}
  >{\raggedright\arraybackslash}p{(\linewidth - 2\tabcolsep) * \real{0.5000}}@{}}
\toprule\noalign{}
\begin{minipage}[b]{\linewidth}\raggedright
字段
\end{minipage} & \begin{minipage}[b]{\linewidth}\raggedright
内容
\end{minipage} \\
\midrule\noalign{}
\endhead
\bottomrule\noalign{}
\endlastfoot
\textbf{主题/材料体系} & 卤化物钙钛矿：MA2PtI6、Cs2AgBiBr6 \\
\textbf{性能属性} & band gap \\
\textbf{构效关系陈述} & 带隙随压力分段线性闭合（两段速率），分段点 ≈ 1.2 GPa 结构转变；二次/指数模型不优于分段线性 \\
\textbf{置信度} & 0.960 \\
\textbf{新颖性} & 0.50 --- 增量新知 \\
\textbf{搜索方法} & bayesian \\
\textbf{Best Score} & 0.000 \\
\textbf{验证状态} & pending \\
\textbf{模型比较} & --- \\
\textbf{外部数据库} & --- \\
\textbf{可提取数据点} & 5 点 / 2 种独立材料 \\
\textbf{证据链} & （证据链条目存在，但需清理格式） \\
\textbf{已知工作} & c9nr07030c 报道 Cs2AgBiBr6 压力带隙演化；1674-1056/adce9e 报道 MA2PtI6 分段闭合 \\
\textbf{增量贡献} & 将两体系的压力响应统一为分段线性标度框架，分段点与结构转变关联 \\
\textbf{描述} & MA2PtI6 带隙在 1.2 GPa 处分两段线性闭合（0.063 → 0.079 eV/GPa，p20）；Cs2AgBiBr6 在 2.3 GPa 相变后红移→蓝移（p15）。假设压力-带隙响应存在分段线性标度，分段点对应结构相变/电子结构重组，且 dEg/dP 与体系压缩性/八面体畸变度相关。用表 C 数据检验线性/分段线性/二次模型的拟合优劣。 \\
\end{longtable}
}

\subsubsection{SPR-PVSK-05 --- ML 联合预测带隙与稳定性：电离能/分解能描述符的独立性检验}\label{spr-pvsk-05-ml-ux8054ux5408ux9884ux6d4bux5e26ux9699ux4e0eux7a33ux5b9aux6027ux7535ux79bbux80fdux5206ux89e3ux80fdux63cfux8ff0ux7b26ux7684ux72ecux7acbux6027ux68c0ux9a8c}

{\def\LTcaptype{none} % do not increment counter
\begin{longtable}[]{@{}
  >{\raggedright\arraybackslash}p{(\linewidth - 2\tabcolsep) * \real{0.5000}}
  >{\raggedright\arraybackslash}p{(\linewidth - 2\tabcolsep) * \real{0.5000}}@{}}
\toprule\noalign{}
\begin{minipage}[b]{\linewidth}\raggedright
字段
\end{minipage} & \begin{minipage}[b]{\linewidth}\raggedright
内容
\end{minipage} \\
\midrule\noalign{}
\endhead
\bottomrule\noalign{}
\endlastfoot
\textbf{主题/材料体系} & 卤化物钙钛矿：ABX3、A2BB'X6 \\
\textbf{性能属性} & band gap \\
\textbf{构效关系陈述} & 描述符空间中带隙与分解能存在正交子空间（可分离）；电离能/容差因子主要预测稳定性，带隙由电子构型描述符主导 \\
\textbf{置信度} & 0.882 \\
\textbf{新颖性} & 0.60 --- 增量新知 \\
\textbf{搜索方法} & bayesian \\
\textbf{Best Score} & 0.000 \\
\textbf{验证状态} & pending \\
\textbf{模型比较} & --- \\
\textbf{外部数据库} & --- \\
\textbf{可提取数据点} & 0 点 / 495 种独立材料 \\
\textbf{证据链} & （证据链条目存在，但需清理格式） \\
\textbf{已知工作} & Landini et al.~ML 筛选双钙钛矿 (p64)；solener.2021.09.030 联合预测 (p19) \\
\textbf{增量贡献} & 从'ML 能预测'推进到'带隙与稳定性在描述符空间可分离'的结构性结论 \\
\textbf{描述} & p60 表明 ML 可同时预测带隙（RMSE 21 meV）与形成能（39 meV/atom）；p65 用 4 类物理描述符（packing/bonding/polarization/electronic identity）筛选 1221 个双钙钛矿；p50 提出电离能作为稳定性判据。假设带隙与稳定性（分解能）在描述符空间具有可分离性------即存在对带隙敏感但稳定性无关（或反之）的独立描述符，使双目标联合优化成为可能。 \\
\end{longtable}
}

\subsubsection{SPR-TE-01 --- 纳米晶尺寸调控 Si-Ge-P 超饱和固溶体的晶格热导率与 ZT}\label{spr-te-01-ux7eb3ux7c73ux6676ux5c3aux5bf8ux8c03ux63a7-si-ge-p-ux8d85ux9971ux548cux56faux6eb6ux4f53ux7684ux6676ux683cux70edux5bfcux7387ux4e0e-zt}

{\def\LTcaptype{none} % do not increment counter
\begin{longtable}[]{@{}
  >{\raggedright\arraybackslash}p{(\linewidth - 2\tabcolsep) * \real{0.5000}}
  >{\raggedright\arraybackslash}p{(\linewidth - 2\tabcolsep) * \real{0.5000}}@{}}
\toprule\noalign{}
\begin{minipage}[b]{\linewidth}\raggedright
字段
\end{minipage} & \begin{minipage}[b]{\linewidth}\raggedright
内容
\end{minipage} \\
\midrule\noalign{}
\endhead
\bottomrule\noalign{}
\endlastfoot
\textbf{主题/材料体系} & 热电材料：Si-Ge-P 超饱和固溶体、Si80Ge20P10Fe1 \\
\textbf{性能属性} & 热电优值 ZT（1000 K） \\
\textbf{构效关系陈述} & 当晶粒尺寸从 50 nm 降至约 10 nm 时，晶格热导率近似随晶粒尺寸倒数增大而下降；在 10 nm 附近，功率因子保持较高水平，ZT 达到最大值。进一步减小晶粒至 5 nm 以下，晶界散射导致载流子迁移率显著下降，功率因子衰减，ZT 下降。预计可在实验中获得 ZT 高于 2.5，并验证 ML 预测的高性能窗口。 \\
\textbf{置信度} & 0.872 \\
\textbf{新颖性} & 0.80 --- 新知（高新颖性） \\
\textbf{搜索方法} & bayesian \\
\textbf{Best Score} & 0.000 \\
\textbf{验证状态} & inconclusive \\
\textbf{模型比较} & candidate\_possible, R²=0.285, p=0.157 \\
\textbf{外部数据库} & materials\_project, oqmd, nomad \\
\textbf{可提取数据点} & 0 点 / 0 种独立材料 \\
\textbf{证据链} & TE002, TE057, TE009 \\
\textbf{已知工作} & 已有文献依据(evidence\_chain 编号: TE002, TE057, TE009)，具体结论需人工/LLM 补写 \\
\textbf{增量贡献} & 相对已确立结论的具体增量待补写(本假设的 expected\_relationship 即拟验证的新规律) \\
\textbf{描述} & 针对 ML 预测高 ZT 材料缺乏实验验证的断层，提出以超饱和 Si-Ge-P 固溶体为模型体系，系统改变球磨时间与低温烧结条件，获得晶粒尺寸在 5--50 nm 的系列样品。假设纳米晶界主要散射中长波声子，而 Fe/P 共掺杂调整电子结构，使功率因子得以保留；因此存在最优晶粒尺寸使晶格热导率大幅下降而载流子迁移率未严重恶化。该假设可通过变温 Hall 效应、电导率/Seebeck 系数和激光闪射法热导率联合测试直接验证，从而在实验中逼近 ML 预测的 ZT 值。 \\
\end{longtable}
}

\subsubsection{SPR-TE-02 --- 卤素掺杂位点与 n 型 SnSe 载流子浓度及 ZT 的关联}\label{spr-te-02-ux5364ux7d20ux63baux6742ux4f4dux70b9ux4e0e-n-ux578b-snse-ux8f7dux6d41ux5b50ux6d53ux5ea6ux53ca-zt-ux7684ux5173ux8054}

{\def\LTcaptype{none} % do not increment counter
\begin{longtable}[]{@{}
  >{\raggedright\arraybackslash}p{(\linewidth - 2\tabcolsep) * \real{0.5000}}
  >{\raggedright\arraybackslash}p{(\linewidth - 2\tabcolsep) * \real{0.5000}}@{}}
\toprule\noalign{}
\begin{minipage}[b]{\linewidth}\raggedright
字段
\end{minipage} & \begin{minipage}[b]{\linewidth}\raggedright
内容
\end{minipage} \\
\midrule\noalign{}
\endhead
\bottomrule\noalign{}
\endlastfoot
\textbf{主题/材料体系} & 热电材料：SnSe、SnSe1-xBrx、SnSe1-xClx、SnSe1-xIx \\
\textbf{性能属性} & n 型载流子浓度与热电优值 ZT \\
\textbf{构效关系陈述} & 随卤素掺杂浓度增加，电子浓度先上升后因缺陷补偿或第二相析出而饱和；Seebeck 系数的绝对值随载流子浓度近似按 n\^{}-1/3 下降，因此功率因子在 n≈5×10\textsuperscript{18--10}19 cm\^{}-3 区间出现峰值。Br 具有适中的离子半径和较低的缺陷形成能，预期比 Cl/I 更能有效提升 n 型 SnSe 的 ZT，最高可达 1.8--2.2。 \\
\textbf{置信度} & 0.850 \\
\textbf{新颖性} & 0.68 --- 新知（中等新颖性） \\
\textbf{搜索方法} & bayesian \\
\textbf{Best Score} & 0.000 \\
\textbf{验证状态} & pending \\
\textbf{模型比较} & candidate\_possible, R²=0.262, p=0.139 \\
\textbf{外部数据库} & --- \\
\textbf{可提取数据点} & 0 点 / 0 种独立材料 \\
\textbf{证据链} & TE122, TE115, TE113 \\
\textbf{已知工作} & 已有文献依据(evidence\_chain 编号: TE122, TE115, TE113)，具体结论需人工/LLM 补写 \\
\textbf{增量贡献} & 相对已确立结论的具体增量待补写(本假设的 expected\_relationship 即拟验证的新规律) \\
\textbf{描述} & 针对 n 型 SnSe 性能远落后 p 型 SnSe 的根因问题，提出以卤素（Cl/Br/I）取代 Se 位并结合 Sn 空位补偿来调控 n 型载流子浓度。假设掺杂剂的离子半径和缺陷形成能决定其固溶度及电离效率，进而控制电子浓度与晶格畸变程度。该假设可通过第一性原理缺陷化学计算预测各掺杂体系的缺陷形成能，再通过实验合成不同卤素掺杂浓度的 SnSe 多晶，结合变温 Hall 和输运测量检验。 \\
\end{longtable}
}

\subsubsection{SPR-TE-03 --- 最优掺杂浓度随温度上移抑制双极效应的实验验证}\label{spr-te-03-ux6700ux4f18ux63baux6742ux6d53ux5ea6ux968fux6e29ux5ea6ux4e0aux79fbux6291ux5236ux53ccux6781ux6548ux5e94ux7684ux5b9eux9a8cux9a8cux8bc1}

{\def\LTcaptype{none} % do not increment counter
\begin{longtable}[]{@{}
  >{\raggedright\arraybackslash}p{(\linewidth - 2\tabcolsep) * \real{0.5000}}
  >{\raggedright\arraybackslash}p{(\linewidth - 2\tabcolsep) * \real{0.5000}}@{}}
\toprule\noalign{}
\begin{minipage}[b]{\linewidth}\raggedright
字段
\end{minipage} & \begin{minipage}[b]{\linewidth}\raggedright
内容
\end{minipage} \\
\midrule\noalign{}
\endhead
\bottomrule\noalign{}
\endlastfoot
\textbf{主题/材料体系} & 热电材料：Bi2Te3、Bi0.5Sb1.5Te3、PbTe \\
\textbf{性能属性} & 高温热电优值 ZT（600 K） \\
\textbf{构效关系陈述} & 与室温最优掺杂浓度相比，600 K 下的最优载流子浓度高约 10--30\%。当掺杂浓度按 n*(T) 随温度升高而增大时，600 K 附近的功率因子衰减被抑制，双极热导率降低，因此 ZT 比固定室温最优掺杂样品提高 15--30\%。这可为宽温域热电材料的分段掺杂设计提供直接实验依据。 \\
\textbf{置信度} & 0.833 \\
\textbf{新颖性} & 0.74 --- 新知（中等新颖性） \\
\textbf{搜索方法} & bayesian \\
\textbf{Best Score} & 0.000 \\
\textbf{验证状态} & pending \\
\textbf{模型比较} & candidate\_possible, R²=0.285, p=0.157 \\
\textbf{外部数据库} & --- \\
\textbf{可提取数据点} & 0 点 / 0 种独立材料 \\
\textbf{证据链} & TE082, TE085, TE090 \\
\textbf{已知工作} & 已有文献依据(evidence\_chain 编号: TE082, TE085, TE090, 预测值 10--30\%（待实验验证，建议方案：元素分析或TGA定量组成变化）)，具体结论需人工/LLM 补写 \\
\textbf{增量贡献} & 相对已确立结论的具体增量待补写(本假设的 expected\_relationship 即拟验证的新规律) \\
\textbf{描述} & 针对双极效应下最优掺杂随温度演化的实验验证缺失，提出在 Bi2Te3 基和 PbTe 材料中，通过制备一系列不同掺杂浓度的样品，系统测量 300--700 K 的电输运和热输运性质，绘制最优载流子浓度 n*(T) 相图。假设随着温度升高，本征激发增强，双极热导率和双极 Seebeck 抵消效应加剧，因此为了抑制双极效应，最优掺杂浓度应向更高值移动。 \\
\end{longtable}
}

\textbf{数值验证结果} - \textbf{预测值（待实验验证）}: \texttt{10–30\%} --- 在 4 篇论文摘要中检索，未找到直接支撑。 建议验证方案: 预测值 10--30\%（待实验验证，建议方案：元素分析或TGA定量组成变化） - **预测值（待实验验证） \textbar{}

\subsubsection{SPR-TE-04 --- 共振掺杂浓度与 DOS 局部异常及 Seebeck 增强的定量标度}\label{spr-te-04-ux5171ux632fux63baux6742ux6d53ux5ea6ux4e0e-dos-ux5c40ux90e8ux5f02ux5e38ux53ca-seebeck-ux589eux5f3aux7684ux5b9aux91cfux6807ux5ea6}

{\def\LTcaptype{none} % do not increment counter
\begin{longtable}[]{@{}
  >{\raggedright\arraybackslash}p{(\linewidth - 2\tabcolsep) * \real{0.5000}}
  >{\raggedright\arraybackslash}p{(\linewidth - 2\tabcolsep) * \real{0.5000}}@{}}
\toprule\noalign{}
\begin{minipage}[b]{\linewidth}\raggedright
字段
\end{minipage} & \begin{minipage}[b]{\linewidth}\raggedright
内容
\end{minipage} \\
\midrule\noalign{}
\endhead
\bottomrule\noalign{}
\endlastfoot
\textbf{主题/材料体系} & 热电材料：HfZrCoSnSb、HfZrCoSnSb:Al、PbTe:K/Na \\
\textbf{性能属性} & Seebeck 系数与功率因子 \\
\textbf{构效关系陈述} & 在 Al 浓度为 0--0.8 at\% 范围内，费米能级处 DOS 异常幅度随浓度近似线性增大，Seebeck 系数随之增加；功率因子在 0.5--0.8 at\% 附近达到峰值。当 Al 浓度超过约 1 at\% 时，体系过渡为金属性杂质带行为，Seebeck 系数快速下降。该标度律可推广到其他共振掺杂体系以替代盲目试错。 \\
\textbf{置信度} & 0.680 \\
\textbf{新颖性} & 0.82 --- 新知（高新颖性） \\
\textbf{搜索方法} & bayesian \\
\textbf{Best Score} & 0.000 \\
\textbf{验证状态} & pending \\
\textbf{模型比较} & --- \\
\textbf{外部数据库} & --- \\
\textbf{可提取数据点} & 0 点 / 0 种独立材料 \\
\textbf{证据链} & TE146, TE087, TE101 \\
\textbf{已知工作} & 已有文献依据(evidence\_chain 编号: TE146, TE087, TE101)，具体结论需人工/LLM 补写 \\
\textbf{增量贡献} & 相对已确立结论的具体增量待补写(本假设的 expected\_relationship 即拟验证的新规律) \\
\textbf{描述} & 针对共振掺杂缺乏定量标度规律的问题，提出在 half-Heusler HfZrCoSnSb 中，通过控制 Al 共振掺杂浓度（0--1.5 at\%）并利用第一性原理计算费米能级附近态密度（DOS）的局部畸变幅度，再与实验测量的 Seebeck 系数和功率因子关联。假设共振杂质产生的 DOS 局部异常幅度越大，Seebeck 系数的增强越显著，但超过临界浓度后杂质带扩展会破坏共振态，导致 Seebeck 骤降。 \\
\end{longtable}
}

\subsubsection{SPR-TE-05 --- Yb 填充方钴矿填充分数对晶格热导率与迁移率的定量影响及批次稳定性}\label{spr-te-05-yb-ux586bux5145ux65b9ux94b4ux77ffux586bux5145ux5206ux6570ux5bf9ux6676ux683cux70edux5bfcux7387ux4e0eux8fc1ux79fbux7387ux7684ux5b9aux91cfux5f71ux54cdux53caux6279ux6b21ux7a33ux5b9aux6027}

{\def\LTcaptype{none} % do not increment counter
\begin{longtable}[]{@{}
  >{\raggedright\arraybackslash}p{(\linewidth - 2\tabcolsep) * \real{0.5000}}
  >{\raggedright\arraybackslash}p{(\linewidth - 2\tabcolsep) * \real{0.5000}}@{}}
\toprule\noalign{}
\begin{minipage}[b]{\linewidth}\raggedright
字段
\end{minipage} & \begin{minipage}[b]{\linewidth}\raggedright
内容
\end{minipage} \\
\midrule\noalign{}
\endhead
\bottomrule\noalign{}
\endlastfoot
\textbf{主题/材料体系} & 热电材料：YbyFe4Sb12、YbyCo4Sb12、CeyFe4Sb12 \\
\textbf{性能属性} & 晶格热导率与 ZT 的批次稳定性 \\
\textbf{构效关系陈述} & 当实际填充分数 y 从 0 增至约 0.25 时，晶格热导率因填充原子的 rattling 散射单调下降约 40--60\%；继续增加填充至 0.4 时，载流子迁移率因缺陷散射增强而明显下降，功率因子受损。因此 ZT 在 y≈0.20--0.25 出现宽峰，且该区域对 y 的敏感度最低，从而解释文献中 0.67--0.89 的 ZT 波动并给出提升批次复现性的合成窗口。 \\
\textbf{置信度} & 0.786 \\
\textbf{新颖性} & 0.65 --- 新知（中等新颖性） \\
\textbf{搜索方法} & bayesian \\
\textbf{Best Score} & 0.000 \\
\textbf{验证状态} & inconclusive \\
\textbf{模型比较} & candidate\_possible, R²=0.285, p=0.157 \\
\textbf{外部数据库} & materials\_project, oqmd, nomad \\
\textbf{可提取数据点} & 0 点 / 0 种独立材料 \\
\textbf{证据链} & TE048, TE035, TE037, TE051 \\
\textbf{已知工作} & 已有文献依据(evidence\_chain 编号: TE048, TE035, TE037, TE051)，具体结论需人工/LLM 补写 \\
\textbf{增量贡献} & 相对已确立结论的具体增量待补写(本假设的 expected\_relationship 即拟验证的新规律) \\
\textbf{描述} & 针对填充方钴矿批次复现性差的问题，提出名义填充分数 y 与实际填充分数之间的偏差是导致 ZT 不一致的主要原因。假设 Yb 填充原子在方钴矿笼子中产生``rattling''效应，可定量降低晶格热导率；但过量填充会引入附加的电子散射并改变载流子浓度。通过制备一系列 y=0--0.4 的 Yb\_yFe4Sb12/Yb\_yCo4Sb12 样品，测定实际填充分数、晶格热导率、载流子迁移率和 ZT，可确定性能最优且对 y 不敏感的填充范围。 \\
\end{longtable}
}

\textbf{数值验证结果} - \textbf{预测值（待实验验证）}: \texttt{40–60\%} --- 在 4 篇论文摘要中检索，未找到直接支撑。 建议验证方案: 预测值 40--60\%（ \textbar{}

\subsubsection{SPR-NMC-01 --- 单晶NMC811中Ni氧化态异质性调控位错辅助裂纹萌生}\label{spr-nmc-01-ux5355ux6676nmc811ux4e2dniux6c27ux5316ux6001ux5f02ux8d28ux6027ux8c03ux63a7ux4f4dux9519ux8f85ux52a9ux88c2ux7eb9ux840cux751f}

{\def\LTcaptype{none} % do not increment counter
\begin{longtable}[]{@{}
  >{\raggedright\arraybackslash}p{(\linewidth - 2\tabcolsep) * \real{0.5000}}
  >{\raggedright\arraybackslash}p{(\linewidth - 2\tabcolsep) * \real{0.5000}}@{}}
\toprule\noalign{}
\begin{minipage}[b]{\linewidth}\raggedright
字段
\end{minipage} & \begin{minipage}[b]{\linewidth}\raggedright
内容
\end{minipage} \\
\midrule\noalign{}
\endhead
\bottomrule\noalign{}
\endlastfoot
\textbf{主题/材料体系} & 高镍正极：LiNi0.8Mn0.1Co0.1O2 (SC-NMC811) \\
\textbf{性能属性} & 裂纹萌生临界应力 \\
\textbf{构效关系陈述} & Ni氧化态空间异质性程度（如Ni价态方差或局部Ni4+分数）与临界应力和裂纹起始位置直接相关：异质性越强，临界应力越低，且裂纹优先从价态突变界面处萌生。 \\
\textbf{置信度} & 0.720 \\
\textbf{新颖性} & 0.88 --- 新知（高新颖性） \\
\textbf{搜索方法} & bayesian \\
\textbf{Best Score} & 0.000 \\
\textbf{验证状态} & inconclusive \\
\textbf{模型比较} & --- \\
\textbf{外部数据库} & materials\_project, oqmd, nomad \\
\textbf{可提取数据点} & 0 点 / 0 种独立材料 \\
\textbf{证据链} & p24, p25 \\
\textbf{已知工作} & 已有文献依据(evidence\_chain 编号: p24, p25)，具体结论需人工/LLM 补写 \\
\textbf{增量贡献} & 相对已确立结论的具体增量待补写(本假设的 expected\_relationship 即拟验证的新规律) \\
\textbf{描述} & 在单晶LiNi0.8Mn0.1Co0.1O2正极颗粒中，Ni氧化态空间异质性（Ni2+/Ni3+/Ni4+分布）会产生局部晶格应变和缺陷能波动，直接影响位错滑移的临界剪切应力。假设同一单晶颗粒中Ni氧化态异质性较大的区域（特别是Ni4+富集或Ni2+还原区域）优先成为裂纹萌生点，导致颗粒整体断裂强度降低。可通过先对同一颗粒进行光谱叠层成像获得Ni价态图，再进行原位微力学压缩实验验证裂纹起始位置与价态异质性的空间关联。 \\
\end{longtable}
}

\subsubsection{SPR-NMC-02 --- 裂纹电解质润湿通过表面重构与氧释放加速容量衰减}\label{spr-nmc-02-ux88c2ux7eb9ux7535ux89e3ux8d28ux6da6ux6e7fux901aux8fc7ux8868ux9762ux91cdux6784ux4e0eux6c27ux91caux653eux52a0ux901fux5bb9ux91cfux8870ux51cf}

{\def\LTcaptype{none} % do not increment counter
\begin{longtable}[]{@{}
  >{\raggedright\arraybackslash}p{(\linewidth - 2\tabcolsep) * \real{0.5000}}
  >{\raggedright\arraybackslash}p{(\linewidth - 2\tabcolsep) * \real{0.5000}}@{}}
\toprule\noalign{}
\begin{minipage}[b]{\linewidth}\raggedright
字段
\end{minipage} & \begin{minipage}[b]{\linewidth}\raggedright
内容
\end{minipage} \\
\midrule\noalign{}
\endhead
\bottomrule\noalign{}
\endlastfoot
\textbf{主题/材料体系} & 高镍正极：LiNi0.8Mn0.1Co0.1O2、LiNiO2、碳酸酯电解液 \\
\textbf{性能属性} & 循环容量保持率 \\
\textbf{构效关系陈述} & 容量保持率随循环圈数下降的速率与裂纹暴露面积×表面重构动力学（电压/温度相关）成正比；在高压（\textgreater4.2 V vs graphite）下，裂纹润湿与氧释放的耦合使衰减加速超出两者单独贡献之和。 \\
\textbf{置信度} & 0.700 \\
\textbf{新颖性} & 0.90 --- 新知（高新颖性） \\
\textbf{搜索方法} & bayesian \\
\textbf{Best Score} & 0.000 \\
\textbf{验证状态} & pending \\
\textbf{模型比较} & --- \\
\textbf{外部数据库} & --- \\
\textbf{可提取数据点} & 0 点 / 0 种独立材料 \\
\textbf{证据链} & p27, p31, p32 \\
\textbf{已知工作} & 已有文献依据(evidence\_chain 编号: p27, p31, p32)，具体结论需人工/LLM 补写 \\
\textbf{增量贡献} & 相对已确立结论的具体增量待补写(本假设的 expected\_relationship 即拟验证的新规律) \\
\textbf{描述} & 在高镍多晶NMC正极中，循环中形成的晶间裂纹被电解质润湿后暴露新的表面，这些新表面在高电压下按LiNiO2表面相图发生重构并释放氧，进一步增加界面阻抗和活性锂损失。假设容量衰减速率与裂纹表面积（或电解质可接触新表面面积）以及充电电压呈正相关，且存在耦合增益：裂纹+润湿比单一裂纹或单一表面重构造成的衰减更快。可通过对比在相同机械损伤下`湿'裂纹与`干'裂纹（如使用惰性气氛或固态电解质阻断润湿）的容量保持率来验证。 \\
\end{longtable}
}

\subsubsection{SPR-NMC-03 --- 高镍层状氧化物中Ni含量与容量保持率呈非单调权衡}\label{spr-nmc-03-ux9ad8ux954dux5c42ux72b6ux6c27ux5316ux7269ux4e2dniux542bux91cfux4e0eux5bb9ux91cfux4fddux6301ux7387ux5448ux975eux5355ux8c03ux6743ux8861}

{\def\LTcaptype{none} % do not increment counter
\begin{longtable}[]{@{}
  >{\raggedright\arraybackslash}p{(\linewidth - 2\tabcolsep) * \real{0.5000}}
  >{\raggedright\arraybackslash}p{(\linewidth - 2\tabcolsep) * \real{0.5000}}@{}}
\toprule\noalign{}
\begin{minipage}[b]{\linewidth}\raggedright
字段
\end{minipage} & \begin{minipage}[b]{\linewidth}\raggedright
内容
\end{minipage} \\
\midrule\noalign{}
\endhead
\bottomrule\noalign{}
\endlastfoot
\textbf{主题/材料体系} & 高镍正极：LiNi0.6Mn0.2Co0.2O2、LiNi0.7Mn0.1Co0.2O2、LiNi0.8Mn0.1Co0.1O2、LiNi0.9Mn0.05Co0.05O2、LiNiO2 \\
\textbf{性能属性} & 循环容量保持效率 (\%) \\
\textbf{构效关系陈述} & 放电容量（0.1C, 2.5-4.3V）随Ni含量线性增加，而300圈循环保持率随Ni含量分段下降：x≤0.8下降平缓，x\textgreater0.9出现加速下降，形成明确拐点；最佳折衷在x≈0.8-0.9区间。 \\
\textbf{置信度} & 0.680 \\
\textbf{新颖性} & 0.82 --- 新知（高新颖性） \\
\textbf{搜索方法} & bayesian \\
\textbf{Best Score} & 0.000 \\
\textbf{验证状态} & validated \\
\textbf{模型比较} & insufficient, R²=0.100 \\
\textbf{外部数据库} & materials\_project, oqmd, nomad \\
\textbf{可提取数据点} & 0 点 / 0 种独立材料 \\
\textbf{证据链} & p22, p24, p32 \\
\textbf{已知工作} & 已有文献依据(evidence\_chain 编号: p22, p24, p32)，具体结论需人工/LLM 补写 \\
\textbf{增量贡献} & 相对已确立结论的具体增量待补写(本假设的 expected\_relationship 即拟验证的新规律) \\
\textbf{描述} & 在LiNi\_x Mn\_y Co\_z O2体系中，随着Ni含量x从0.6增加到1.0，首次放电容量增加，但由于Ni4+的强氧化性和表面/体相不稳定性导致循环保持率下降。假设在x≈0.8-0.9之间存在容量-保持率帕累托最优窗口，而x\textgreater0.9后保持率急剧恶化。该假设可通过合成一系列x=0.6/0.7/0.8/0.9/1.0的单晶或球形多晶样品，在相同电压窗口和电解液下进行标准化循环测试加以检验。 \\
\end{longtable}
}

\subsubsection{SPR-NMC-04 --- 涂层材料的Li离子电导率与界面副反应阻抗之间的帕累托最优设计}\label{spr-nmc-04-ux6d82ux5c42ux6750ux6599ux7684liux79bbux5b50ux7535ux5bfcux7387ux4e0eux754cux9762ux526fux53cdux5e94ux963bux6297ux4e4bux95f4ux7684ux5e15ux7d2fux6258ux6700ux4f18ux8bbeux8ba1}

{\def\LTcaptype{none} % do not increment counter
\begin{longtable}[]{@{}
  >{\raggedright\arraybackslash}p{(\linewidth - 2\tabcolsep) * \real{0.5000}}
  >{\raggedright\arraybackslash}p{(\linewidth - 2\tabcolsep) * \real{0.5000}}@{}}
\toprule\noalign{}
\begin{minipage}[b]{\linewidth}\raggedright
字段
\end{minipage} & \begin{minipage}[b]{\linewidth}\raggedright
内容
\end{minipage} \\
\midrule\noalign{}
\endhead
\bottomrule\noalign{}
\endlastfoot
\textbf{主题/材料体系} & 高镍正极：α-AlF3、α-Al2O3、m-ZrO2、c-MgO、SiO2、LiNi0.8Mn0.1Co0.1O2 \\
\textbf{性能属性} & 倍率容量保持效率 (\%) \\
\textbf{构效关系陈述} & log(倍率容量保持率)随涂层厚度线性下降，下降斜率与log(D\_Li)成反比；副反应抑制率随厚度单调上升；在交叉点处存在最佳厚度h* = f(D\_Li, 界面反应速率常数)。 \\
\textbf{置信度} & 0.650 \\
\textbf{新颖性} & 0.78 --- 新知（中等新颖性） \\
\textbf{搜索方法} & bayesian \\
\textbf{Best Score} & 0.000 \\
\textbf{验证状态} & pending \\
\textbf{模型比较} & --- \\
\textbf{外部数据库} & --- \\
\textbf{可提取数据点} & 0 点 / 0 种独立材料 \\
\textbf{证据链} & p39, p42 \\
\textbf{已知工作} & 已有文献依据(evidence\_chain 编号: p39, p42)，具体结论需人工/LLM 补写 \\
\textbf{增量贡献} & 相对已确立结论的具体增量待补写(本假设的 expected\_relationship 即拟验证的新规律) \\
\textbf{描述} & 正极表面涂层（Al2O3、AlF3、ZrO2、MgO、SiO2等）对Li离子传输的阻碍与对电解液副反应的阻挡构成矛盾。假设涂层材料的体积Li扩散系数D\_Li和电子绝缘性/热力学稳定性共同决定最优厚度：存在一个临界涂层厚度 h\emph{，低于h}不能有效抑制副反应，高于h\emph{则导致倍率性能快速下降；且D\_Li越高，h}可越厚而不损失倍率性能。可通过在同一高镍正极上控制不同涂层材料与厚度并进行倍率-循环测试验证。 \\
\end{longtable}
}

\subsubsection{SPR-NMC-05 --- Ni氧化态升高降低阳离子混排迁移势垒，促进层状结构向岩盐相降解}\label{spr-nmc-05-niux6c27ux5316ux6001ux5347ux9ad8ux964dux4f4eux9633ux79bbux5b50ux6df7ux6392ux8fc1ux79fbux52bfux5792ux4fc3ux8fdbux5c42ux72b6ux7ed3ux6784ux5411ux5ca9ux76d0ux76f8ux964dux89e3}

{\def\LTcaptype{none} % do not increment counter
\begin{longtable}[]{@{}
  >{\raggedright\arraybackslash}p{(\linewidth - 2\tabcolsep) * \real{0.5000}}
  >{\raggedright\arraybackslash}p{(\linewidth - 2\tabcolsep) * \real{0.5000}}@{}}
\toprule\noalign{}
\begin{minipage}[b]{\linewidth}\raggedright
字段
\end{minipage} & \begin{minipage}[b]{\linewidth}\raggedright
内容
\end{minipage} \\
\midrule\noalign{}
\endhead
\bottomrule\noalign{}
\endlastfoot
\textbf{主题/材料体系} & 高镍正极：LiNi0.8Mn0.1Co0.1O2、LiNiO2 \\
\textbf{性能属性} & 阳离子混排度 / 容量衰减速率 \\
\textbf{构效关系陈述} & Ni平均价态每增加0.1（对应Li脱出量增加），Li层Ni占据分数呈指数增长，且容量衰减速率与混排度呈线性正相关；表面重构层厚度与Ni4+浓度同步增加。 \\
\textbf{置信度} & 0.630 \\
\textbf{新颖性} & 0.85 --- 新知（高新颖性） \\
\textbf{搜索方法} & bayesian \\
\textbf{Best Score} & 0.000 \\
\textbf{验证状态} & pending \\
\textbf{模型比较} & --- \\
\textbf{外部数据库} & --- \\
\textbf{可提取数据点} & 0 点 / 0 种独立材料 \\
\textbf{证据链} & p27, p32, p24 \\
\textbf{已知工作} & 已有文献依据(evidence\_chain 编号: p27, p32, p24)，具体结论需人工/LLM 补写 \\
\textbf{增量贡献} & 相对已确立结论的具体增量待补写(本假设的 expected\_relationship 即拟验证的新规律) \\
\textbf{描述} & 根据氧化诱导阳离子无序理论，高Ni3+/Ni4+含量使Ni经四面体中间体迁移至Li层。假设镍氧化态越高（如充电态SOC高），阳离子迁移势垒越低，阳离子混排度越高，导致电压和容量衰减。可以通过原位XAS/STEM定量不同电位下Ni价态和Li/Ni交换度，并测量对应的容量衰减速率。 \\
\end{longtable}
}

\subsubsection{SPR-NMC-06 --- 水洗引入的质子残留通过近表面Li/H交换加速高镍正极容量衰减与阻抗增长}\label{spr-nmc-06-ux6c34ux6d17ux5f15ux5165ux7684ux8d28ux5b50ux6b8bux7559ux901aux8fc7ux8fd1ux8868ux9762lihux4ea4ux6362ux52a0ux901fux9ad8ux954dux6b63ux6781ux5bb9ux91cfux8870ux51cfux4e0eux963bux6297ux589eux957f}

{\def\LTcaptype{none} % do not increment counter
\begin{longtable}[]{@{}
  >{\raggedright\arraybackslash}p{(\linewidth - 2\tabcolsep) * \real{0.5000}}
  >{\raggedright\arraybackslash}p{(\linewidth - 2\tabcolsep) * \real{0.5000}}@{}}
\toprule\noalign{}
\begin{minipage}[b]{\linewidth}\raggedright
字段
\end{minipage} & \begin{minipage}[b]{\linewidth}\raggedright
内容
\end{minipage} \\
\midrule\noalign{}
\endhead
\bottomrule\noalign{}
\endlastfoot
\textbf{主题/材料体系} & 高镍正极：LiNi₀.₈₃Co₀.₁₂Mn₀.₀₅O₂ (NCM831205) \\
\textbf{性能属性} & 质子诱导容量衰减速率 \\
\textbf{构效关系陈述} & 质子含量（ppm 级）每增加一个数量级，100 圈容量衰减速率与界面阻抗呈近似线性上升（低含量段即显著），衰减增量与 Ni 含量正相关；质子残留与 Zr 涂层（p46）、过锂化（p44）存在交互效应。 \\
\textbf{置信度} & 0.650 \\
\textbf{新颖性} & 0.85 --- 新知（高新颖性） \\
\textbf{搜索方法} & bayesian \\
\textbf{Best Score} & 0.000 \\
\textbf{验证状态} & validated \\
\textbf{模型比较} & --- \\
\textbf{外部数据库} & materials\_project, oqmd, nomad \\
\textbf{可提取数据点} & 0 点 / 0 种独立材料 \\
\textbf{证据链} & p48, p46, p44 \\
\textbf{已知工作} & 已有文献依据(evidence\_chain 编号: p48, p46, p44)，具体结论需人工/LLM 补写 \\
\textbf{增量贡献} & 相对已确立结论的具体增量待补写(本假设的 expected\_relationship 即拟验证的新规律) \\
\textbf{描述} & 高镍NCM正极水洗除杂时近表面插层锂离子与水中质子发生交换（Li+/H+交换），质子残留在充放电中脱出并参与界面副反应。p48揭示即使低质子含量也会引发大性能变化：容量衰减与阻抗随质子含量上升。假设质子含量（水洗程度）与容量衰减速率、界面阻抗呈单调正相关，且存在阈值效应------低质子含量即显著影响。可通过系统制备不同水洗程度（质子含量梯度）样品，用OEMS+EIS定量质子含量-阻抗-衰减速率关系验证，并考察质子与涂层/过锂化工艺的交互。 \\
\end{longtable}
}

\subsubsection{SPR-SE-01 --- 卤素比例调控卤化物固态电解质室温电导率与湿度稳定性的非线性关系}\label{spr-se-01-ux5364ux7d20ux6bd4ux4f8bux8c03ux63a7ux5364ux5316ux7269ux56faux6001ux7535ux89e3ux8d28ux5ba4ux6e29ux7535ux5bfcux7387ux4e0eux6e7fux5ea6ux7a33ux5b9aux6027ux7684ux975eux7ebfux6027ux5173ux7cfb}

{\def\LTcaptype{none} % do not increment counter
\begin{longtable}[]{@{}
  >{\raggedright\arraybackslash}p{(\linewidth - 2\tabcolsep) * \real{0.5000}}
  >{\raggedright\arraybackslash}p{(\linewidth - 2\tabcolsep) * \real{0.5000}}@{}}
\toprule\noalign{}
\begin{minipage}[b]{\linewidth}\raggedright
字段
\end{minipage} & \begin{minipage}[b]{\linewidth}\raggedright
内容
\end{minipage} \\
\midrule\noalign{}
\endhead
\bottomrule\noalign{}
\endlastfoot
\textbf{主题/材料体系} & 固态电解质：Li3YCl6、Li3YCl6-xBrx、Li3InCl6、Li2ZrCl6 \\
\textbf{性能属性} & 室温离子电导率与湿度稳定性（H2O暴露后电导率保持率） \\
\textbf{构效关系陈述} & 预期室温离子电导率与Br取代量x呈火山型关系，在x≈1.0--1.5达到峰值；湿度稳定性随x增加单调下降。最优x位于电导率增益与稳定性损失交叉点附近，该点可同时获得\textgreater1 mS/cm的室温电导率和可接受的湿度稳定性。 \\
\textbf{置信度} & 0.720 \\
\textbf{新颖性} & 0.70 --- 新知（中等新颖性） \\
\textbf{搜索方法} & bayesian \\
\textbf{Best Score} & 0.000 \\
\textbf{验证状态} & pending \\
\textbf{模型比较} & --- \\
\textbf{外部数据库} & --- \\
\textbf{可提取数据点} & 0 点 / 0 种独立材料 \\
\textbf{证据链} & P040, P070, P006 \\
\textbf{已知工作} & 已有文献依据(evidence\_chain 编号: P040, P070, P006)，具体结论需人工/LLM 补写 \\
\textbf{增量贡献} & 相对已确立结论的具体增量待补写(本假设的 expected\_relationship 即拟验证的新规律) \\
\textbf{描述} & 在Li3YCl6等卤化物固态电解质中，用Br-或I-部分取代Cl-可以扩张晶格体积、增大锂离子迁移通道尺寸，从而降低迁移活化能、提升室温离子电导率。然而，较大且极化率较高的卤素离子同时会增强吸湿性和与水分的反应活性，导致湿度稳定性下降。因此，室温离子电导率随Br/I取代比例呈火山型变化，而湿度稳定性呈单调递减，二者之间存在最优卤素配比窗口。本假说通过系统合成Li3YCl6-xBrx系列并测量电导率与湿度暴露后的结构/电导保持率来验证。 \\
\end{longtable}
}

\subsubsection{SPR-SE-02 --- 氧/钼双掺杂对硫银锗矿硫化物电解质电导率与H2S释放量的双目标优化窗口}\label{spr-se-02-ux6c27ux94bcux53ccux63baux6742ux5bf9ux786bux94f6ux9517ux77ffux786bux5316ux7269ux7535ux89e3ux8d28ux7535ux5bfcux7387ux4e0eh2sux91caux653eux91cfux7684ux53ccux76eeux6807ux4f18ux5316ux7a97ux53e3}

{\def\LTcaptype{none} % do not increment counter
\begin{longtable}[]{@{}
  >{\raggedright\arraybackslash}p{(\linewidth - 2\tabcolsep) * \real{0.5000}}
  >{\raggedright\arraybackslash}p{(\linewidth - 2\tabcolsep) * \real{0.5000}}@{}}
\toprule\noalign{}
\begin{minipage}[b]{\linewidth}\raggedright
字段
\end{minipage} & \begin{minipage}[b]{\linewidth}\raggedright
内容
\end{minipage} \\
\midrule\noalign{}
\endhead
\bottomrule\noalign{}
\endlastfoot
\textbf{主题/材料体系} & 固态电解质：Li6PS5Cl、Li6PS5Cl1-xOx、Mo-O共掺杂Li6PS5Cl、Li6PS5I \\
\textbf{性能属性} & 室温离子电导率及H2S生成速率 \\
\textbf{构效关系陈述} & 预期存在一个掺杂浓度区间（例如O含量0.1--0.3），使室温电导率保持≥1 mS/cm的同时，H2S释放量降低50\%以上。氧/钼双掺杂的效果优于单一氧掺杂，因为Mo可进一步稳定晶格并促进Li+跃迁。 \\
\textbf{置信度} & 0.786 \\
\textbf{新颖性} & 0.75 --- 新知（中等新颖性） \\
\textbf{搜索方法} & bayesian \\
\textbf{Best Score} & 0.000 \\
\textbf{验证状态} & pending \\
\textbf{模型比较} & --- \\
\textbf{外部数据库} & --- \\
\textbf{可提取数据点} & 0 点 / 0 种独立材料 \\
\textbf{证据链} & P024, P031, P027, P023 \\
\textbf{已知工作} & 已有文献依据(evidence\_chain 编号: P024, P031, P027, P023)，具体结论需人工/LLM 补写 \\
\textbf{增量贡献} & 相对已确立结论的具体增量待补写(本假设的 expected\_relationship 即拟验证的新规律) \\
\textbf{描述} & 硫化物固态电解质在潮湿环境中易水解产生H2S，而掺杂氧或金属氧化物（如MoO2）可能同时改善电导率和湿度稳定性。本假说认为，氧掺杂引入更强的局部共价键和更致密的阴离子框架，一方面通过缺陷工程维持甚至提升Li+迁移率，另一方面氧位点与水分子的结合能高于硫位点，从而抑制水解反应。通过系统调控Li6PS5Cl中的O含量及Mo-O共掺杂比例，可以突破传统``电导率-湿度稳定性''此消彼长的权衡关系。 \\
\end{longtable}
}

\textbf{数值验证结果} - \textbf{预测值（待实验验证）}: \texttt{50\%} --- 在 5 篇论文摘要中检索，未找到直接支撑。 建议验证方案: 预测值 50\%（待实验验证，建议方案：元素分析或TGA定量组成变化） \textbar{}

\subsubsection{SPR-SE-03 --- 高熵阳离子无序度与固态电解质室温离子电导率的火山型关系及机器学习预测闭环}\label{spr-se-03-ux9ad8ux71b5ux9633ux79bbux5b50ux65e0ux5e8fux5ea6ux4e0eux56faux6001ux7535ux89e3ux8d28ux5ba4ux6e29ux79bbux5b50ux7535ux5bfcux7387ux7684ux706bux5c71ux578bux5173ux7cfbux53caux673aux5668ux5b66ux4e60ux9884ux6d4bux95edux73af}

{\def\LTcaptype{none} % do not increment counter
\begin{longtable}[]{@{}
  >{\raggedright\arraybackslash}p{(\linewidth - 2\tabcolsep) * \real{0.5000}}
  >{\raggedright\arraybackslash}p{(\linewidth - 2\tabcolsep) * \real{0.5000}}@{}}
\toprule\noalign{}
\begin{minipage}[b]{\linewidth}\raggedright
字段
\end{minipage} & \begin{minipage}[b]{\linewidth}\raggedright
内容
\end{minipage} \\
\midrule\noalign{}
\endhead
\bottomrule\noalign{}
\endlastfoot
\textbf{主题/材料体系} & 固态电解质：Li3(Y0.2In0.2Sc0.2Ho0.2Er0.2)Cl6、Li3YCl6、Li3InCl6、高熵硫化物Li6PS5Cl基固溶体 \\
\textbf{性能属性} & 室温离子电导率 \\
\textbf{构效关系陈述} & 预期室温离子电导率随构型熵S\_cfg增大先升高后降低，呈火山型曲线，最优S\_cfg对应最大电导率。机器学习模型若以局域配位环境为特征，可定量捕捉该非线性关系；实验验证后，模型预测误差将用于指导下一轮候选合成，形成闭环。 \\
\textbf{置信度} & 0.650 \\
\textbf{新颖性} & 0.85 --- 新知（高新颖性） \\
\textbf{搜索方法} & bayesian \\
\textbf{Best Score} & 0.000 \\
\textbf{验证状态} & pending \\
\textbf{模型比较} & --- \\
\textbf{外部数据库} & --- \\
\textbf{可提取数据点} & 0 点 / 0 种独立材料 \\
\textbf{证据链} & P068, P080, P079, P074 \\
\textbf{已知工作} & 已有文献依据(evidence\_chain 编号: P068, P080, P079, P074)，具体结论需人工/LLM 补写 \\
\textbf{增量贡献} & 相对已确立结论的具体增量待补写(本假设的 expected\_relationship 即拟验证的新规律) \\
\textbf{描述} & 机器学习筛选出的高熵固态电解质候选材料缺乏实验验证，且高熵带来的晶格畸变对离子输运的影响尚不明确。本假说提出，在卤化物/硫化物框架中引入多种阳离子（如Y, In, Sc, Ho, Er）形成高熵固溶体，构型熵增加会拓宽锂离子迁移势垒分布，但适度无序可增加迁移通道的连通性和瓶颈尺寸；过高熵值则产生严重晶格畸变，堵塞迁移路径。利用OBELiX等数据库训练机器学习模型，预测不同高熵组成的电导率，并用实验测量值校验，可建立构型熵-电导率的定量关系并改进模型外推能力。 \\
\end{longtable}
}

\subsubsection{SPR-SE-04 --- Lewis酸性基团浓度对聚合物电解质锂离子迁移数与离子电导率权衡的非单调调控}\label{spr-se-04-lewisux9178ux6027ux57faux56e2ux6d53ux5ea6ux5bf9ux805aux5408ux7269ux7535ux89e3ux8d28ux9502ux79bbux5b50ux8fc1ux79fbux6570ux4e0eux79bbux5b50ux7535ux5bfcux7387ux6743ux8861ux7684ux975eux5355ux8c03ux8c03ux63a7}

{\def\LTcaptype{none} % do not increment counter
\begin{longtable}[]{@{}
  >{\raggedright\arraybackslash}p{(\linewidth - 2\tabcolsep) * \real{0.5000}}
  >{\raggedright\arraybackslash}p{(\linewidth - 2\tabcolsep) * \real{0.5000}}@{}}
\toprule\noalign{}
\begin{minipage}[b]{\linewidth}\raggedright
字段
\end{minipage} & \begin{minipage}[b]{\linewidth}\raggedright
内容
\end{minipage} \\
\midrule\noalign{}
\endhead
\bottomrule\noalign{}
\endlastfoot
\textbf{主题/材料体系} & 固态电解质：PEO-LiTFSI、含硼酸酯（Lewis酸）PEO基聚合物、含TFSAM配位性阴离子的聚合物电解质 \\
\textbf{性能属性} & 锂离子迁移数（t+）与离子电导率（σ） \\
\textbf{构效关系陈述} & 预期随着Lewis酸基团浓度增加，t+从\textasciitilde0.2单调上升至0.5以上，而σ先升后降。原因是适量Lewis酸增强盐解离并促进Li+传输，但过量Lewis酸会提高玻璃化转变温度或形成交联网络，降低链段运动能力。因此t+-σ权衡曲线存在一个最优Lewis酸浓度，使t+与σ同时优于未改性PEO。 \\
\textbf{置信度} & 0.620 \\
\textbf{新颖性} & 0.80 --- 新知（高新颖性） \\
\textbf{搜索方法} & bayesian \\
\textbf{Best Score} & 0.000 \\
\textbf{验证状态} & pending \\
\textbf{模型比较} & --- \\
\textbf{外部数据库} & --- \\
\textbf{可提取数据点} & 0 点 / 0 种独立材料 \\
\textbf{证据链} & P038, P039, P059, P001 \\
\textbf{已知工作} & 已有文献依据(evidence\_chain 编号: P038, P039, P059, P001)，具体结论需人工/LLM 补写 \\
\textbf{增量贡献} & 相对已确立结论的具体增量待补写(本假设的 expected\_relationship 即拟验证的新规律) \\
\textbf{描述} & 传统PEO基聚合物电解质中，锂离子迁移数t+低（\textasciitilde0.2）且与离子电导率σ存在权衡。弱配位阴离子（如TFSI）并不能解耦Li+与链段运动，而引入Lewis酸性基团（如硼酸酯）或配位性阴离子（如TFSAM）可能通过可逆配位阴离子来降低阴离子迁移率，从而提升t+。本假说系统研究Lewis酸含量及阴离子类型对t+和σ的影响，以期找到协同优化窗口。 \\
\end{longtable}
}

\subsubsection{SPR-SE-05 --- 亲锂成核层+电子阻挡层双层界面设计对固态电解质界面电阻的普适降低效应}\label{spr-se-05-ux4eb2ux9502ux6210ux6838ux5c42ux7535ux5b50ux963bux6321ux5c42ux53ccux5c42ux754cux9762ux8bbeux8ba1ux5bf9ux56faux6001ux7535ux89e3ux8d28ux754cux9762ux7535ux963bux7684ux666eux9002ux964dux4f4eux6548ux5e94}

{\def\LTcaptype{none} % do not increment counter
\begin{longtable}[]{@{}
  >{\raggedright\arraybackslash}p{(\linewidth - 2\tabcolsep) * \real{0.5000}}
  >{\raggedright\arraybackslash}p{(\linewidth - 2\tabcolsep) * \real{0.5000}}@{}}
\toprule\noalign{}
\begin{minipage}[b]{\linewidth}\raggedright
字段
\end{minipage} & \begin{minipage}[b]{\linewidth}\raggedright
内容
\end{minipage} \\
\midrule\noalign{}
\endhead
\bottomrule\noalign{}
\endlastfoot
\textbf{主题/材料体系} & 固态电解质：Li6PS5Cl、Li7La3Zr2O12 (LLZO)、Li3YCl6、Ag涂层、LiF涂层、Ag@COOH-CNTs复合层 \\
\textbf{性能属性} & Li/SSE界面电阻（Ω cm²） \\
\textbf{构效关系陈述} & 预期在多种SSE上，Ag亲锂层负责诱导均匀锂沉积、降低成核过电位，LiF电子阻挡层负责抑制电子跨界面传输、减少副反应，二者协同可将界面电阻从\textgreater100 Ω cm²降至\textless10 Ω cm²，且该效应与SSE阴离子种类无关。界面电阻降低幅度还可能与SSE的还原稳定性相关，即越不稳定的SSE受益越大。 \\
\textbf{置信度} & 0.660 \\
\textbf{新颖性} & 0.75 --- 新知（中等新颖性） \\
\textbf{搜索方法} & bayesian \\
\textbf{Best Score} & 0.000 \\
\textbf{验证状态} & pending \\
\textbf{模型比较} & --- \\
\textbf{外部数据库} & --- \\
\textbf{可提取数据点} & 0 点 / 0 种独立材料 \\
\textbf{证据链} & P047, P017, P060, P058, P046 \\
\textbf{已知工作} & 已有文献依据(evidence\_chain 编号: P047, P017, P060, P058)，具体结论需人工/LLM 补写 \\
\textbf{增量贡献} & 相对已确立结论的具体增量待补写(本假设的 expected\_relationship 即拟验证的新规律) \\
\textbf{描述} & 实验室中通过Ag@COOH-CNTs、LiF富集层、LiCux三维网络等界面工程可将Li/SSE界面电阻降至0.25 Ω cm²，但工程级全固态电池仍受\textgreater100 Ω cm²界面电阻困扰。本假说提出，这些策略成功的共性在于同时满足``亲锂成核位点''和``电子阻挡层''两个功能；若能将其提炼为普适的双层界面设计（如Ag/LiF双层），则可跨硫化物、氧化物、卤化物SSE实现界面电阻的大幅降低。 \\
\end{longtable}
}

\subsection{3. 统计附表}\label{ux7edfux8ba1ux9644ux8868}

\subsubsection{3.1 按主题统计}\label{ux6309ux4e3bux9898ux7edfux8ba1}

{\def\LTcaptype{none} % do not increment counter
\begin{longtable}[]{@{}
  >{\raggedright\arraybackslash}p{(\linewidth - 14\tabcolsep) * \real{0.0822}}
  >{\raggedright\arraybackslash}p{(\linewidth - 14\tabcolsep) * \real{0.1096}}
  >{\raggedright\arraybackslash}p{(\linewidth - 14\tabcolsep) * \real{0.1096}}
  >{\raggedright\arraybackslash}p{(\linewidth - 14\tabcolsep) * \real{0.1096}}
  >{\raggedright\arraybackslash}p{(\linewidth - 14\tabcolsep) * \real{0.0822}}
  >{\raggedright\arraybackslash}p{(\linewidth - 14\tabcolsep) * \real{0.1507}}
  >{\raggedright\arraybackslash}p{(\linewidth - 14\tabcolsep) * \real{0.1507}}
  >{\raggedright\arraybackslash}p{(\linewidth - 14\tabcolsep) * \real{0.2055}}@{}}
\toprule\noalign{}
\begin{minipage}[b]{\linewidth}\raggedright
主题
\end{minipage} & \begin{minipage}[b]{\linewidth}\raggedright
假设数
\end{minipage} & \begin{minipage}[b]{\linewidth}\raggedright
已验证
\end{minipage} & \begin{minipage}[b]{\linewidth}\raggedright
待验证
\end{minipage} & \begin{minipage}[b]{\linewidth}\raggedright
待定
\end{minipage} & \begin{minipage}[b]{\linewidth}\raggedright
平均置信度
\end{minipage} & \begin{minipage}[b]{\linewidth}\raggedright
平均新颖性
\end{minipage} & \begin{minipage}[b]{\linewidth}\raggedright
外部数据库覆盖
\end{minipage} \\
\midrule\noalign{}
\endhead
\bottomrule\noalign{}
\endlastfoot
MOF CO₂ 捕获 & 5 & 1 & 4 & 0 & 0.790 & 0.85 & materials\_project \\
MOF CO₂ 捕获（e2e v4 重跑） & 5 & 2 & 3 & 0 & 0.859 & 0.75 & hmof\_core\_mof, materials\_project, nomad, oqmd \\
卤化物钙钛矿 & 5 & 2 & 3 & 0 & 0.947 & 0.58 & materials\_project, nomad, oqmd \\
热电材料 & 5 & 0 & 3 & 2 & 0.804 & 0.74 & materials\_project, nomad, oqmd \\
高镍正极 & 6 & 2 & 3 & 1 & 0.672 & 0.85 & materials\_project, nomad, oqmd \\
固态电解质 & 5 & 0 & 5 & 0 & 0.687 & 0.77 & --- \\
\end{longtable}
}

\subsubsection{3.2 搜索方法分布}\label{ux641cux7d22ux65b9ux6cd5ux5206ux5e03}

\begin{itemize}
\tightlist
\item
  \textbf{bayesian}：29 条
\item
  \textbf{hybrid}：1 条
\item
  \textbf{mcts}：1 条
\end{itemize}

\subsubsection{3.3 验证状态分布}\label{ux9a8cux8bc1ux72b6ux6001ux5206ux5e03}

\begin{itemize}
\tightlist
\item
  \textbf{validated}：7 条
\item
  \textbf{pending}：21 条
\item
  \textbf{inconclusive}：3 条
\end{itemize}

\begin{center}\rule{0.5\linewidth}{0.5pt}\end{center}

\emph{本文档由 \texttt{scripts/build\_route\_a\_docs.py} 自动生成（2026 年 8 月），数据源为 6 个主题的 \texttt{discovery/hypotheses.json}。}

\end{document}
